% ============================================================
% Documentación de Endpoints API - Backend FastAPI
% Generado a partir del código fuente en /backend/routes, /backend/schemas y /backend/models
% ============================================================

\subsection{Autenticación}

\begin{enumerate}
    \item 1. http://”URLPROYECTO”.cl/auth/login
    \begin{enumerate}
        \item a. Tipo de petición:
            \begin{itemize}
               \item i. “post”
            \end{itemize}
        \item b. Parámetros a ingresar, el tipo de datos y restricciones
            \begin{enumerate}
                \item i. username (body, x-www-form-urlencoded)
                    \begin{enumerate}
                        \item 1. String
                        \item 2. Requerido
                        \item 3. Corresponde al email del empleado
                    \end{enumerate}
                \item ii. password
                    \begin{enumerate}
                        \item 1. String
                        \item 2. Requerido
                        \item 3. Se verifica con hash Argon2
                    \end{enumerate}
            \end{enumerate}
        \item c. Respuesta: 200 \{access\_token, token\_type: "bearer"\}; 401 contraseña incorrecta; 403 usuario desactivado; 404 email no encontrado.
    \end{enumerate}
\end{enumerate}

\subsection{Tipos de Estado (/tipos-estado)}

\begin{enumerate}
    \item 1. http://”URLPROYECTO”.cl/tipos-estado/
    \begin{enumerate}
        \item a. Tipo de petición:
            \begin{itemize}
               \item i. “get”
            \end{itemize}
        \item b. Parámetros a ingresar, el tipo de datos y restricciones
            \begin{enumerate}
                \item i. Sin parámetros de entrada (requiere token de autenticación)
            \end{enumerate}
    \end{enumerate}

    \item 2. http://”URLPROYECTO”.cl/tipos-estado/\{id\_tipoe\}
    \begin{enumerate}
        \item a. Tipo de petición:
            \begin{itemize}
               \item i. “get”
            \end{itemize}
        \item b. Parámetros a ingresar, el tipo de datos y restricciones
            \begin{enumerate}
                \item i. id\_tipoe (parámetro de ruta)
                    \begin{enumerate}
                        \item 1. Number (entero)
                        \item 2. Requerido
                    \end{enumerate}
            \end{enumerate}
    \end{enumerate}

    \item 3. http://”URLPROYECTO”.cl/tipos-estado/
    \begin{enumerate}
        \item a. Tipo de petición:
            \begin{itemize}
               \item i. “post”
            \end{itemize}
        \item b. Parámetros a ingresar, el tipo de datos y restricciones
            \begin{enumerate}
                \item i. nombre
                    \begin{enumerate}
                        \item 1. String
                        \item 2. Requerido
                    \end{enumerate}
                \item ii. descr
                    \begin{enumerate}
                        \item 1. String
                        \item 2. No requerido - Default null
                    \end{enumerate}
                \item iii. operatividad
                    \begin{enumerate}
                        \item 1. Boolean
                        \item 2. No requerido - Default true
                    \end{enumerate}
            \end{enumerate}
    \end{enumerate}

    \item 4. http://”URLPROYECTO”.cl/tipos-estado/\{id\_tipoe\}
    \begin{enumerate}
        \item a. Tipo de petición:
            \begin{itemize}
               \item i. “put”
            \end{itemize}
        \item b. Parámetros a ingresar, el tipo de datos y restricciones
            \begin{enumerate}
                \item i. id\_tipoe (parámetro de ruta)
                    \begin{enumerate}
                        \item 1. Number (entero)
                        \item 2. Requerido
                    \end{enumerate}
                \item ii. nombre
                    \begin{enumerate}
                        \item 1. String
                        \item 2. No requerido - Default null
                    \end{enumerate}
                \item iii. descr
                    \begin{enumerate}
                        \item 1. String
                        \item 2. No requerido - Default null
                    \end{enumerate}
                \item iv. operatividad
                    \begin{enumerate}
                        \item 1. Boolean
                        \item 2. No requerido - Default null
                    \end{enumerate}
            \end{enumerate}
    \end{enumerate}

    \item 5. http://”URLPROYECTO”.cl/tipos-estado/\{id\_tipoe\}
    \begin{enumerate}
        \item a. Tipo de petición:
            \begin{itemize}
               \item i. “delete”
            \end{itemize}
        \item b. Parámetros a ingresar, el tipo de datos y restricciones
            \begin{enumerate}
                \item i. id\_tipoe (parámetro de ruta)
                    \begin{enumerate}
                        \item 1. Number (entero)
                        \item 2. Requerido
                    \end{enumerate}
            \end{enumerate}
        \item c. Nota: borrado lógico (operatividad = false)
    \end{enumerate}
\end{enumerate}

\subsection{Departamentos (/departamento)}

\begin{enumerate}
    \item 1. http://”URLPROYECTO”.cl/departamento/
    \begin{enumerate}
        \item a. Tipo de petición:
            \begin{itemize}
               \item i. “get”
            \end{itemize}
        \item b. Parámetros a ingresar, el tipo de datos y restricciones
            \begin{enumerate}
                \item i. Sin parámetros de entrada
            \end{enumerate}
    \end{enumerate}

    \item 2. http://”URLPROYECTO”.cl/departamento/\{id\_dpto\}
    \begin{enumerate}
        \item a. Tipo de petición:
            \begin{itemize}
               \item i. “get”
            \end{itemize}
        \item b. Parámetros a ingresar, el tipo de datos y restricciones
            \begin{enumerate}
                \item i. id\_dpto (parámetro de ruta)
                    \begin{enumerate}
                        \item 1. Number (entero)
                        \item 2. Requerido
                    \end{enumerate}
            \end{enumerate}
    \end{enumerate}

    \item 3. http://”URLPROYECTO”.cl/departamento/
    \begin{enumerate}
        \item a. Tipo de petición:
            \begin{itemize}
               \item i. “post”
            \end{itemize}
        \item b. Parámetros a ingresar, el tipo de datos y restricciones
            \begin{enumerate}
                \item i. id\_org
                    \begin{enumerate}
                        \item 1. Number (entero)
                        \item 2. No requerido - Default null
                        \item 3. Debe existir en la tabla Empresa
                    \end{enumerate}
                \item ii. zona
                    \begin{enumerate}
                        \item 1. String (enum)
                        \item 2. Requerido
                        \item 3. Valores permitidos: "Administracion", "Operaciones", "Produccion"
                    \end{enumerate}
                \item iii. nom\_dpto
                    \begin{enumerate}
                        \item 1. String
                        \item 2. Requerido
                        \item 3. Tamaño máximo de 100 caracteres
                        \item 4. No puede repetirse (case-insensitive) dentro de la misma empresa
                    \end{enumerate}
            \end{enumerate}
        \item c. Nota: este endpoint no exige autenticación (posible pendiente de seguridad)
    \end{enumerate}

    \item 4. http://”URLPROYECTO”.cl/departamento/\{id\_dpto\}
    \begin{enumerate}
        \item a. Tipo de petición:
            \begin{itemize}
               \item i. “put”
            \end{itemize}
        \item b. Parámetros a ingresar, el tipo de datos y restricciones
            \begin{enumerate}
                \item i. id\_dpto (parámetro de ruta)
                    \begin{enumerate}
                        \item 1. Number (entero)
                        \item 2. Requerido
                    \end{enumerate}
                \item ii. id\_org
                    \begin{enumerate}
                        \item 1. Number (entero)
                        \item 2. No requerido - Default null
                    \end{enumerate}
                \item iii. zona
                    \begin{enumerate}
                        \item 1. String (enum)
                        \item 2. No requerido - Default null
                        \item 3. Valores permitidos: "Administracion", "Operaciones", "Produccion"
                    \end{enumerate}
                \item iv. nom\_dpto
                    \begin{enumerate}
                        \item 1. String
                        \item 2. No requerido - Default null
                        \item 3. Tamaño máximo de 100 caracteres
                    \end{enumerate}
            \end{enumerate}
    \end{enumerate}
\end{enumerate}

\subsection{Empleados (/empleado)}

\begin{enumerate}
    \item 1. http://”URLPROYECTO”.cl/empleado/
    \begin{enumerate}
        \item a. Tipo de petición:
            \begin{itemize}
               \item i. “get”
            \end{itemize}
        \item b. Parámetros a ingresar, el tipo de datos y restricciones
            \begin{enumerate}
                \item i. Sin parámetros de entrada
            \end{enumerate}
    \end{enumerate}

    \item 2. http://”URLPROYECTO”.cl/empleado/\{id\_emp\}
    \begin{enumerate}
        \item a. Tipo de petición:
            \begin{itemize}
               \item i. “get”
            \end{itemize}
        \item b. Parámetros a ingresar, el tipo de datos y restricciones
            \begin{enumerate}
                \item i. id\_emp (parámetro de ruta)
                    \begin{enumerate}
                        \item 1. Number (entero)
                        \item 2. Requerido
                    \end{enumerate}
            \end{enumerate}
    \end{enumerate}

    \item 3. http://”URLPROYECTO”.cl/empleado/
    \begin{enumerate}
        \item a. Tipo de petición:
            \begin{itemize}
               \item i. “post”
            \end{itemize}
        \item b. Parámetros a ingresar, el tipo de datos y restricciones
            \begin{enumerate}
                \item i. rut
                    \begin{enumerate}
                        \item 1. String
                        \item 2. Requerido
                        \item 3. Único en el sistema
                        \item 4. Tamaño máximo de 50 caracteres
                    \end{enumerate}
                \item ii. nombre
                    \begin{enumerate}
                        \item 1. String
                        \item 2. Requerido
                        \item 3. Tamaño máximo de 50 caracteres
                    \end{enumerate}
                \item iii. email
                    \begin{enumerate}
                        \item 1. String (formato email)
                        \item 2. Requerido
                        \item 3. Único en el sistema
                        \item 4. Tamaño máximo de 50 caracteres
                    \end{enumerate}
                \item iv. cargo
                    \begin{enumerate}
                        \item 1. String
                        \item 2. Requerido
                        \item 3. Tamaño máximo de 100 caracteres
                        \item 4. Si no existe en el catálogo Cargo se crea automáticamente
                    \end{enumerate}
                \item v. planta
                    \begin{enumerate}
                        \item 1. String
                        \item 2. Requerido
                        \item 3. Tamaño máximo de 100 caracteres
                        \item 4. Si no existe en el catálogo Planta se crea automáticamente
                    \end{enumerate}
                \item vi. nom\_dpto
                    \begin{enumerate}
                        \item 1. String
                        \item 2. No requerido - Default null
                        \item 3. Tamaño máximo de 100 caracteres
                        \item 4. Debe existir previamente en Departamento
                    \end{enumerate}
                \item vii. operatividad
                    \begin{enumerate}
                        \item 1. Boolean
                        \item 2. No requerido - Default true
                    \end{enumerate}
                \item viii. id\_sub
                    \begin{enumerate}
                        \item 1. Number (entero)
                        \item 2. No requerido - Default null
                        \item 3. Debe existir en Sub-área
                    \end{enumerate}
            \end{enumerate}
        \item c. Nota: este endpoint no exige autenticación (posible pendiente de seguridad)
    \end{enumerate}

    \item 4. http://”URLPROYECTO”.cl/empleado/\{id\_emp\}
    \begin{enumerate}
        \item a. Tipo de petición:
            \begin{itemize}
               \item i. “put”
            \end{itemize}
        \item b. Parámetros a ingresar, el tipo de datos y restricciones
            \begin{enumerate}
                \item i. id\_emp (parámetro de ruta)
                    \begin{enumerate}
                        \item 1. Number (entero)
                        \item 2. Requerido
                    \end{enumerate}
                \item ii. rut, nombre, email, cargo, planta, nom\_dpto, operatividad, id\_sub
                    \begin{enumerate}
                        \item 1. Mismos tipos y restricciones que en la creación
                        \item 2. No requeridos - Default null (actualización parcial)
                    \end{enumerate}
                \item iii. id\_dpto
                    \begin{enumerate}
                        \item 1. Number (entero)
                        \item 2. No requerido - Default null
                        \item 3. Si se envía junto con id\_sub, la subárea debe pertenecer a ese departamento
                    \end{enumerate}
            \end{enumerate}
    \end{enumerate}

    \item 5. http://”URLPROYECTO”.cl/empleado/\{id\_emp\}
    \begin{enumerate}
        \item a. Tipo de petición:
            \begin{itemize}
               \item i. “delete”
            \end{itemize}
        \item b. Parámetros a ingresar, el tipo de datos y restricciones
            \begin{enumerate}
                \item i. id\_emp (parámetro de ruta)
                    \begin{enumerate}
                        \item 1. Number (entero)
                        \item 2. Requerido
                    \end{enumerate}
            \end{enumerate}
        \item c. Nota: borrado lógico (operatividad = false)
    \end{enumerate}
\end{enumerate}

\subsection{Usuarios del Sistema (/usuario\_sistema)}

\begin{enumerate}
    \item 1. http://”URLPROYECTO”.cl/usuario\_sistema/
    \begin{enumerate}
        \item a. Tipo de petición:
            \begin{itemize}
               \item i. “get”
            \end{itemize}
        \item b. Parámetros a ingresar, el tipo de datos y restricciones
            \begin{enumerate}
                \item i. Sin parámetros de entrada
            \end{enumerate}
    \end{enumerate}

    \item 2. http://”URLPROYECTO”.cl/usuario\_sistema/\{id\_emp\}
    \begin{enumerate}
        \item a. Tipo de petición:
            \begin{itemize}
               \item i. “get”
            \end{itemize}
        \item b. Parámetros a ingresar, el tipo de datos y restricciones
            \begin{enumerate}
                \item i. id\_emp (parámetro de ruta)
                    \begin{enumerate}
                        \item 1. Number (entero)
                        \item 2. Requerido
                    \end{enumerate}
            \end{enumerate}
    \end{enumerate}

    \item 3. http://”URLPROYECTO”.cl/usuario\_sistema/
    \begin{enumerate}
        \item a. Tipo de petición:
            \begin{itemize}
               \item i. “post”
            \end{itemize}
        \item b. Parámetros a ingresar, el tipo de datos y restricciones
            \begin{enumerate}
                \item i. id\_emp
                    \begin{enumerate}
                        \item 1. Number (entero)
                        \item 2. Requerido
                        \item 3. Debe existir como Empleado y no tener ya un usuario asociado
                    \end{enumerate}
                \item ii. password
                    \begin{enumerate}
                        \item 1. String
                        \item 2. Requerido
                        \item 3. Tamaño mínimo de 8 caracteres
                        \item 4. Tamaño máximo de 128 caracteres
                    \end{enumerate}
                \item iii. rol
                    \begin{enumerate}
                        \item 1. String (enum)
                        \item 2. Requerido
                        \item 3. Valores permitidos: "editor", "auditor", "admin"
                    \end{enumerate}
                \item iv. vigente
                    \begin{enumerate}
                        \item 1. Boolean
                        \item 2. No requerido - Default true
                    \end{enumerate}
            \end{enumerate}
    \end{enumerate}

    \item 4. http://”URLPROYECTO”.cl/usuario\_sistema/\{id\_emp\}
    \begin{enumerate}
        \item a. Tipo de petición:
            \begin{itemize}
               \item i. “put”
            \end{itemize}
        \item b. Parámetros a ingresar, el tipo de datos y restricciones
            \begin{enumerate}
                \item i. id\_emp (parámetro de ruta)
                    \begin{enumerate}
                        \item 1. Number (entero)
                        \item 2. Requerido
                    \end{enumerate}
                \item ii. password
                    \begin{enumerate}
                        \item 1. String
                        \item 2. No requerido - Default null
                        \item 3. Tamaño mínimo de 8 caracteres
                        \item 4. Tamaño máximo de 128 caracteres
                    \end{enumerate}
                \item iii. rol
                    \begin{enumerate}
                        \item 1. String (enum)
                        \item 2. No requerido - Default null
                        \item 3. Valores permitidos: "editor", "auditor", "admin"
                    \end{enumerate}
                \item iv. vigente
                    \begin{enumerate}
                        \item 1. Boolean
                        \item 2. No requerido - Default null
                        \item 3. No puede establecerse en false sobre el propio usuario autenticado
                    \end{enumerate}
            \end{enumerate}
    \end{enumerate}

    \item 5. http://”URLPROYECTO”.cl/usuario\_sistema/\{id\_emp\}
    \begin{enumerate}
        \item a. Tipo de petición:
            \begin{itemize}
               \item i. “delete”
            \end{itemize}
        \item b. Parámetros a ingresar, el tipo de datos y restricciones
            \begin{enumerate}
                \item i. id\_emp (parámetro de ruta)
                    \begin{enumerate}
                        \item 1. Number (entero)
                        \item 2. Requerido
                        \item 3. No puede ser el propio usuario autenticado
                    \end{enumerate}
            \end{enumerate}
        \item c. Nota: borrado lógico (vigente = false)
    \end{enumerate}
\end{enumerate}

\subsection{Licencias (/licencia)}

\begin{enumerate}
    \item 1. http://”URLPROYECTO”.cl/licencia/
    \begin{enumerate}
        \item a. Tipo de petición:
            \begin{itemize}
               \item i. “get”
            \end{itemize}
        \item b. Parámetros a ingresar, el tipo de datos y restricciones
            \begin{enumerate}
                \item i. Sin parámetros de entrada
            \end{enumerate}
    \end{enumerate}

    \item 2. http://”URLPROYECTO”.cl/licencia/\{id\_lic\}
    \begin{enumerate}
        \item a. Tipo de petición:
            \begin{itemize}
               \item i. “get”
            \end{itemize}
        \item b. Parámetros a ingresar, el tipo de datos y restricciones
            \begin{enumerate}
                \item i. id\_lic (parámetro de ruta)
                    \begin{enumerate}
                        \item 1. Number (entero)
                        \item 2. Requerido
                    \end{enumerate}
            \end{enumerate}
    \end{enumerate}

    \item 3. http://”URLPROYECTO”.cl/licencia/
    \begin{enumerate}
        \item a. Tipo de petición:
            \begin{itemize}
               \item i. “post”
            \end{itemize}
        \item b. Parámetros a ingresar, el tipo de datos y restricciones
            \begin{enumerate}
                \item i. fecha\_ini
                    \begin{enumerate}
                        \item 1. Date
                        \item 2. No requerido - Default null
                    \end{enumerate}
                \item ii. fecha\_exp
                    \begin{enumerate}
                        \item 1. Date
                        \item 2. No requerido - Default null
                    \end{enumerate}
                \item iii. id\_prov
                    \begin{enumerate}
                        \item 1. Number (entero)
                        \item 2. Requerido
                        \item 3. Debe existir en Proveedor de Licencia
                    \end{enumerate}
                \item iv. id\_ver
                    \begin{enumerate}
                        \item 1. Number (entero)
                        \item 2. Requerido
                        \item 3. Debe existir en Versión de Licencia y corresponder al id\_prov indicado
                    \end{enumerate}
            \end{enumerate}
    \end{enumerate}

    \item 4. http://”URLPROYECTO”.cl/licencia/\{id\_lic\}
    \begin{enumerate}
        \item a. Tipo de petición:
            \begin{itemize}
               \item i. “put”
            \end{itemize}
        \item b. Parámetros a ingresar, el tipo de datos y restricciones
            \begin{enumerate}
                \item i. id\_lic (parámetro de ruta)
                    \begin{enumerate}
                        \item 1. Number (entero)
                        \item 2. Requerido
                    \end{enumerate}
                \item ii. fecha\_ini, fecha\_exp, id\_prov, id\_ver
                    \begin{enumerate}
                        \item 1. Mismos tipos y restricciones que en la creación
                        \item 2. No requeridos - Default null
                    \end{enumerate}
            \end{enumerate}
    \end{enumerate}
\end{enumerate}

\subsection{Empleado - Licencia (/emp\_licencias)}

\begin{enumerate}
    \item 1. http://”URLPROYECTO”.cl/emp\_licencias/
    \begin{enumerate}
        \item a. Tipo de petición:
            \begin{itemize}
               \item i. “post”
            \end{itemize}
        \item b. Parámetros a ingresar, el tipo de datos y restricciones
            \begin{enumerate}
                \item i. id\_emp
                    \begin{enumerate}
                        \item 1. Number (entero)
                        \item 2. Requerido
                        \item 3. Debe existir en Empleado
                    \end{enumerate}
                \item ii. id\_lic
                    \begin{enumerate}
                        \item 1. Number (entero)
                        \item 2. Requerido
                        \item 3. Debe existir en Licencia
                        \item 4. La combinación id\_emp + id\_lic no puede estar duplicada
                    \end{enumerate}
            \end{enumerate}
    \end{enumerate}

    \item 2. http://”URLPROYECTO”.cl/emp\_licencias/\{id\_emp\}/\{id\_lic\}
    \begin{enumerate}
        \item a. Tipo de petición:
            \begin{itemize}
               \item i. “delete”
            \end{itemize}
        \item b. Parámetros a ingresar, el tipo de datos y restricciones
            \begin{enumerate}
                \item i. id\_emp (parámetro de ruta)
                    \begin{enumerate}
                        \item 1. Number (entero)
                        \item 2. Requerido
                    \end{enumerate}
                \item ii. id\_lic (parámetro de ruta)
                    \begin{enumerate}
                        \item 1. Number (entero)
                        \item 2. Requerido
                    \end{enumerate}
            \end{enumerate}
        \item c. Nota: borrado físico
    \end{enumerate}
\end{enumerate}

\subsection{Facturas (/factura)}

\begin{enumerate}
    \item 1. http://”URLPROYECTO”.cl/factura/
    \begin{enumerate}
        \item a. Tipo de petición:
            \begin{itemize}
               \item i. “get”
            \end{itemize}
        \item b. Parámetros a ingresar, el tipo de datos y restricciones
            \begin{enumerate}
                \item i. Sin parámetros de entrada
            \end{enumerate}
    \end{enumerate}

    \item 2. http://”URLPROYECTO”.cl/factura/\{id\_fact\}
    \begin{enumerate}
        \item a. Tipo de petición:
            \begin{itemize}
               \item i. “get”
            \end{itemize}
        \item b. Parámetros a ingresar, el tipo de datos y restricciones
            \begin{enumerate}
                \item i. id\_fact (parámetro de ruta)
                    \begin{enumerate}
                        \item 1. Number (entero)
                        \item 2. Requerido
                    \end{enumerate}
            \end{enumerate}
    \end{enumerate}

    \item 3. http://”URLPROYECTO”.cl/factura/
    \begin{enumerate}
        \item a. Tipo de petición:
            \begin{itemize}
               \item i. “post”
            \end{itemize}
        \item b. Parámetros a ingresar, el tipo de datos y restricciones
            \begin{enumerate}
                \item i. id\_prov\_fact
                    \begin{enumerate}
                        \item 1. Number (entero)
                        \item 2. Requerido
                        \item 3. Debe existir en Proveedor de Factura
                    \end{enumerate}
                \item ii. fecha\_em
                    \begin{enumerate}
                        \item 1. Date
                        \item 2. No requerido - Default null
                    \end{enumerate}
                \item iii. num\_folio
                    \begin{enumerate}
                        \item 1. Number (entero)
                        \item 2. Requerido
                    \end{enumerate}
                \item iv. ruta\_pdf\_fact
                    \begin{enumerate}
                        \item 1. String
                        \item 2. Requerido
                        \item 3. Tamaño máximo de 200 caracteres
                    \end{enumerate}
            \end{enumerate}
    \end{enumerate}

    \item 4. http://”URLPROYECTO”.cl/factura/\{id\_fact\}
    \begin{enumerate}
        \item a. Tipo de petición:
            \begin{itemize}
               \item i. “put”
            \end{itemize}
        \item b. Parámetros a ingresar, el tipo de datos y restricciones
            \begin{enumerate}
                \item i. id\_fact (parámetro de ruta)
                    \begin{enumerate}
                        \item 1. Number (entero)
                        \item 2. Requerido
                    \end{enumerate}
                \item ii. id\_prov\_fact, fecha\_em, num\_folio, ruta\_pdf\_fact
                    \begin{enumerate}
                        \item 1. Mismos tipos y restricciones que en la creación
                        \item 2. No requeridos - Default null
                    \end{enumerate}
            \end{enumerate}
    \end{enumerate}
\end{enumerate}

\subsection{Activos TI (/ActivoTi)}

\begin{enumerate}
    \item 1. http://”URLPROYECTO”.cl/ActivoTi/
    \begin{enumerate}
        \item a. Tipo de petición:
            \begin{itemize}
               \item i. “get”
            \end{itemize}
        \item b. Parámetros a ingresar, el tipo de datos y restricciones
            \begin{enumerate}
                \item i. Sin parámetros de entrada
            \end{enumerate}
    \end{enumerate}

    \item 2. http://”URLPROYECTO”.cl/ActivoTi/\{id\_act\}
    \begin{enumerate}
        \item a. Tipo de petición:
            \begin{itemize}
               \item i. “get”
            \end{itemize}
        \item b. Parámetros a ingresar, el tipo de datos y restricciones
            \begin{enumerate}
                \item i. id\_act (parámetro de ruta)
                    \begin{enumerate}
                        \item 1. Number (entero)
                        \item 2. Requerido
                    \end{enumerate}
            \end{enumerate}
    \end{enumerate}

    \item 3. http://”URLPROYECTO”.cl/ActivoTi/
    \begin{enumerate}
        \item a. Tipo de petición:
            \begin{itemize}
               \item i. “post”
            \end{itemize}
        \item b. Parámetros a ingresar, el tipo de datos y restricciones (campos comunes)
            \begin{enumerate}
                \item i. tipo
                    \begin{enumerate}
                        \item 1. String (enum discriminador)
                        \item 2. Requerido
                        \item 3. Valores permitidos: "equipo\_computo", "periferico"
                    \end{enumerate}
                \item ii. id\_modelo
                    \begin{enumerate}
                        \item 1. Number (entero)
                        \item 2. No requerido - Default null
                        \item 3. Debe existir en Modelo
                    \end{enumerate}
                \item iii. num\_serie
                    \begin{enumerate}
                        \item 1. String
                        \item 2. No requerido - Default null
                        \item 3. Único en el sistema
                        \item 4. Tamaño máximo de 50 caracteres
                    \end{enumerate}
                \item iv. origen
                    \begin{enumerate}
                        \item 1. String (enum)
                        \item 2. Requerido
                        \item 3. Valores permitidos: "empresa", "personal"
                    \end{enumerate}
                \item v. fecha\_recep
                    \begin{enumerate}
                        \item 1. Date
                        \item 2. No requerido - Default null
                        \item 3. No puede ser una fecha futura
                    \end{enumerate}
                \item vi. id\_fact
                    \begin{enumerate}
                        \item 1. Number (entero)
                        \item 2. No requerido - Default null
                        \item 3. Debe existir en Factura
                    \end{enumerate}
                \item vii. id\_tipoe
                    \begin{enumerate}
                        \item 1. Number (entero)
                        \item 2. Requerido
                        \item 3. Debe existir en Tipo de Estado
                    \end{enumerate}
                \item viii. id\_dispo
                    \begin{enumerate}
                        \item 1. Number (entero)
                        \item 2. Requerido
                        \item 3. Debe existir en Disponibilidad
                        \item 4. Debe ser compatible con id\_tipoe (ej. un activo "bueno" no puede tener disponibilidad "de\_baja" o "en\_mantencion")
                    \end{enumerate}
            \end{enumerate}
        \item c. Parámetros adicionales si tipo = "equipo\_computo"
            \begin{enumerate}
                \item i. ram\_capac
                    \begin{enumerate}
                        \item 1. Number (entero, enum)
                        \item 2. No requerido - Default null
                        \item 3. Valores permitidos: 2, 4, 6, 8, 10, 12, 16, 20
                    \end{enumerate}
                \item ii. ram\_tipo
                    \begin{enumerate}
                        \item 1. String (enum)
                        \item 2. No requerido - Default null
                        \item 3. Valores permitidos: "DDR3", "DDR4", "DDR5"
                    \end{enumerate}
                \item iii. cpu
                    \begin{enumerate}
                        \item 1. String
                        \item 2. No requerido - Default null
                        \item 3. Tamaño máximo de 50 caracteres
                    \end{enumerate}
                \item iv. mac\_ethernet
                    \begin{enumerate}
                        \item 1. String
                        \item 2. No requerido - Default null
                        \item 3. Debe cumplir formato de dirección MAC (ej. AA:BB:CC:DD:EE:FF)
                    \end{enumerate}
                \item v. mac\_wifi
                    \begin{enumerate}
                        \item 1. String
                        \item 2. No requerido - Default null
                        \item 3. Debe cumplir formato de dirección MAC
                    \end{enumerate}
                \item vi. disco\_tipo
                    \begin{enumerate}
                        \item 1. String (enum)
                        \item 2. No requerido - Default null
                        \item 3. Valores permitidos: "HDD", "SSD", "NVMe", "eMMC"
                    \end{enumerate}
                \item vii. disco\_unidad
                    \begin{enumerate}
                        \item 1. String (enum)
                        \item 2. No requerido - Default null
                        \item 3. Valores permitidos: "MB", "GB", "TB"
                    \end{enumerate}
                \item viii. disco\_capac
                    \begin{enumerate}
                        \item 1. Number (entero)
                        \item 2. No requerido - Default null
                    \end{enumerate}
                \item ix. host\_name
                    \begin{enumerate}
                        \item 1. String
                        \item 2. No requerido - Default null
                        \item 3. Tamaño máximo de 50 caracteres
                    \end{enumerate}
                \item x. id\_tipoeq
                    \begin{enumerate}
                        \item 1. Number (entero)
                        \item 2. No requerido - Default null
                        \item 3. Debe existir en Tipo de Equipo
                    \end{enumerate}
                \item xi. id\_so
                    \begin{enumerate}
                        \item 1. Number (entero)
                        \item 2. No requerido - Default null
                        \item 3. Debe existir en Edición de Sistema Operativo
                    \end{enumerate}
            \end{enumerate}
        \item d. Parámetros adicionales si tipo = "periferico"
            \begin{enumerate}
                \item i. valor\_tec
                    \begin{enumerate}
                        \item 1. Number (decimal, 10 dígitos, 2 decimales)
                        \item 2. No requerido - Default null
                    \end{enumerate}
                \item ii. u\_medida
                    \begin{enumerate}
                        \item 1. String
                        \item 2. No requerido - Default null
                    \end{enumerate}
                \item iii. obs\_perif
                    \begin{enumerate}
                        \item 1. String
                        \item 2. No requerido - Default null
                    \end{enumerate}
                \item iv. id\_tipoper
                    \begin{enumerate}
                        \item 1. Number (entero)
                        \item 2. No requerido - Default null
                        \item 3. Debe existir en Tipo de Periférico
                    \end{enumerate}
            \end{enumerate}
    \end{enumerate}

    \item 4. http://”URLPROYECTO”.cl/ActivoTi/\{id\_act\}
    \begin{enumerate}
        \item a. Tipo de petición:
            \begin{itemize}
               \item i. “put”
            \end{itemize}
        \item b. Parámetros a ingresar, el tipo de datos y restricciones
            \begin{enumerate}
                \item i. id\_act (parámetro de ruta)
                    \begin{enumerate}
                        \item 1. Number (entero)
                        \item 2. Requerido
                    \end{enumerate}
                \item ii. Todos los campos del punto anterior (excepto "tipo", que no puede modificarse)
                    \begin{enumerate}
                        \item 1. Mismos tipos y restricciones que en la creación
                        \item 2. No requeridos - Default null (actualización parcial)
                    \end{enumerate}
            \end{enumerate}
    \end{enumerate}

    \item 5. http://”URLPROYECTO”.cl/ActivoTi/\{id\_act\}
    \begin{enumerate}
        \item a. Tipo de petición:
            \begin{itemize}
               \item i. “delete”
            \end{itemize}
        \item b. Parámetros a ingresar, el tipo de datos y restricciones
            \begin{enumerate}
                \item i. id\_act (parámetro de ruta)
                    \begin{enumerate}
                        \item 1. Number (entero)
                        \item 2. Requerido
                    \end{enumerate}
            \end{enumerate}
        \item c. Nota: implementado como borrado lógico, aunque el campo referenciado en el código no corresponde a una columna real del modelo (hallazgo técnico a revisar)
    \end{enumerate}
\end{enumerate}

\subsection{Historial de Estados (/historial\_estado)}

\begin{enumerate}
    \item 1. http://”URLPROYECTO”.cl/historial\_estado/
    \begin{enumerate}
        \item a. Tipo de petición:
            \begin{itemize}
               \item i. “get”
            \end{itemize}
        \item b. Parámetros a ingresar, el tipo de datos y restricciones
            \begin{enumerate}
                \item i. id\_act (parámetro de query)
                    \begin{enumerate}
                        \item 1. Number (entero)
                        \item 2. No requerido - Default null
                    \end{enumerate}
                \item ii. id\_tipoe (parámetro de query)
                    \begin{enumerate}
                        \item 1. Number (entero)
                        \item 2. No requerido - Default null
                    \end{enumerate}
            \end{enumerate}
    \end{enumerate}

    \item 2. http://”URLPROYECTO”.cl/historial\_estado/\{id\_hes\}
    \begin{enumerate}
        \item a. Tipo de petición:
            \begin{itemize}
               \item i. “get”
            \end{itemize}
        \item b. Parámetros a ingresar, el tipo de datos y restricciones
            \begin{enumerate}
                \item i. id\_hes (parámetro de ruta)
                    \begin{enumerate}
                        \item 1. Number (entero)
                        \item 2. Requerido
                    \end{enumerate}
            \end{enumerate}
    \end{enumerate}

    \item 3. http://”URLPROYECTO”.cl/historial\_estado/
    \begin{enumerate}
        \item a. Tipo de petición:
            \begin{itemize}
               \item i. “post”
            \end{itemize}
        \item b. Parámetros a ingresar, el tipo de datos y restricciones
            \begin{enumerate}
                \item i. descr
                    \begin{enumerate}
                        \item 1. String
                        \item 2. No requerido - Default null
                    \end{enumerate}
                \item ii. causa
                    \begin{enumerate}
                        \item 1. String
                        \item 2. No requerido - Default null
                        \item 3. Tamaño máximo de 50 caracteres
                    \end{enumerate}
                \item iii. id\_tipoe
                    \begin{enumerate}
                        \item 1. Number (entero)
                        \item 2. Requerido
                        \item 3. Debe existir en Tipo de Estado
                    \end{enumerate}
                \item iv. id\_act
                    \begin{enumerate}
                        \item 1. Number (entero)
                        \item 2. Requerido
                        \item 3. Debe existir en Activo TI
                    \end{enumerate}
                \item v. id\_emp\_reg
                    \begin{enumerate}
                        \item 1. Number (entero)
                        \item 2. No requerido - Default null
                        \item 3. Debe existir en Usuario del Sistema
                    \end{enumerate}
            \end{enumerate}
        \item c. Nota: al registrar, actualiza automáticamente ActivoTi.id\_tipoe con el nuevo estado
    \end{enumerate}
\end{enumerate}

\subsection{Asignaciones (/asignaciones)}

\begin{enumerate}
    \item 1. http://”URLPROYECTO”.cl/asignaciones/
    \begin{enumerate}
        \item a. Tipo de petición:
            \begin{itemize}
               \item i. “get”
            \end{itemize}
        \item b. Parámetros a ingresar, el tipo de datos y restricciones
            \begin{enumerate}
                \item i. Sin parámetros de entrada
            \end{enumerate}
    \end{enumerate}

    \item 2. http://”URLPROYECTO”.cl/asignaciones/\{id\_asig\}
    \begin{enumerate}
        \item a. Tipo de petición:
            \begin{itemize}
               \item i. “get”
            \end{itemize}
        \item b. Parámetros a ingresar, el tipo de datos y restricciones
            \begin{enumerate}
                \item i. id\_asig (parámetro de ruta)
                    \begin{enumerate}
                        \item 1. Number (entero)
                        \item 2. Requerido
                    \end{enumerate}
            \end{enumerate}
    \end{enumerate}

    \item 3. http://”URLPROYECTO”.cl/asignaciones/
    \begin{enumerate}
        \item a. Tipo de petición:
            \begin{itemize}
               \item i. “post”
            \end{itemize}
        \item b. Parámetros a ingresar, el tipo de datos y restricciones
            \begin{enumerate}
                \item i. fecha\_entr
                    \begin{enumerate}
                        \item 1. Date
                        \item 2. Requerido
                        \item 3. No puede ser una fecha futura
                    \end{enumerate}
                \item ii. fecha\_dev
                    \begin{enumerate}
                        \item 1. Date
                        \item 2. No requerido - Default null
                        \item 3. No puede ser una fecha futura
                    \end{enumerate}
                \item iii. ruta\_acta\_entr
                    \begin{enumerate}
                        \item 1. String
                        \item 2. No requerido - Default null
                        \item 3. Tamaño máximo de 255 caracteres
                    \end{enumerate}
                \item iv. id\_act
                    \begin{enumerate}
                        \item 1. Number (entero)
                        \item 2. Requerido
                        \item 3. Debe existir en Activo TI y no tener una asignación vigente sin devolución
                    \end{enumerate}
                \item v. id\_emp
                    \begin{enumerate}
                        \item 1. Number (entero)
                        \item 2. Requerido
                        \item 3. Debe existir en Empleado y estar operativo
                    \end{enumerate}
            \end{enumerate}
        \item c. Nota: id\_emp\_reg se asigna automáticamente en el servidor con el usuario autenticado (no se recibe del cliente)
    \end{enumerate}

    \item 4. http://”URLPROYECTO”.cl/asignaciones/\{id\_asig\}
    \begin{enumerate}
        \item a. Tipo de petición:
            \begin{itemize}
               \item i. “put”
            \end{itemize}
        \item b. Parámetros a ingresar, el tipo de datos y restricciones
            \begin{enumerate}
                \item i. id\_asig (parámetro de ruta)
                    \begin{enumerate}
                        \item 1. Number (entero)
                        \item 2. Requerido
                    \end{enumerate}
                \item ii. fecha\_dev
                    \begin{enumerate}
                        \item 1. Date
                        \item 2. Requerido
                        \item 3. No puede ser una fecha futura
                    \end{enumerate}
                \item iii. ruta\_acta\_entr
                    \begin{enumerate}
                        \item 1. String
                        \item 2. No requerido - Default null
                        \item 3. Tamaño máximo de 255 caracteres
                    \end{enumerate}
            \end{enumerate}
        \item c. Nota: solo se puede actualizar si la asignación aún no tiene fecha\_dev registrada
    \end{enumerate}
\end{enumerate}

\subsection{Checklist (/Checklist)}

\begin{enumerate}
    \item 1. http://”URLPROYECTO”.cl/Checklist/
    \begin{enumerate}
        \item a. Tipo de petición:
            \begin{itemize}
               \item i. “get”
            \end{itemize}
        \item b. Parámetros a ingresar, el tipo de datos y restricciones
            \begin{enumerate}
                \item i. Sin parámetros de entrada
            \end{enumerate}
    \end{enumerate}

    \item 2. http://”URLPROYECTO”.cl/Checklist/\{id\_check\}
    \begin{enumerate}
        \item a. Tipo de petición:
            \begin{itemize}
               \item i. “get”
            \end{itemize}
        \item b. Parámetros a ingresar, el tipo de datos y restricciones
            \begin{enumerate}
                \item i. id\_check (parámetro de ruta)
                    \begin{enumerate}
                        \item 1. Number (entero)
                        \item 2. Requerido
                    \end{enumerate}
            \end{enumerate}
    \end{enumerate}

    \item 3. http://”URLPROYECTO”.cl/Checklist/
    \begin{enumerate}
        \item a. Tipo de petición:
            \begin{itemize}
               \item i. “post”
            \end{itemize}
        \item b. Parámetros a ingresar, el tipo de datos y restricciones
            \begin{enumerate}
                \item i. nombre
                    \begin{enumerate}
                        \item 1. String
                        \item 2. Requerido
                        \item 3. Tamaño máximo de 50 caracteres
                    \end{enumerate}
                \item ii. fecha\_ini
                    \begin{enumerate}
                        \item 1. Date
                        \item 2. No requerido - Default null
                    \end{enumerate}
                \item iii. fecha\_ter
                    \begin{enumerate}
                        \item 1. Date
                        \item 2. No requerido - Default null
                    \end{enumerate}
                \item iv. id\_emp
                    \begin{enumerate}
                        \item 1. Number (entero)
                        \item 2. No requerido - Default null
                        \item 3. Debe existir en Usuario del Sistema
                    \end{enumerate}
            \end{enumerate}
    \end{enumerate}

    \item 4. http://”URLPROYECTO”.cl/Checklist/\{id\_check\}
    \begin{enumerate}
        \item a. Tipo de petición:
            \begin{itemize}
               \item i. “put”
            \end{itemize}
        \item b. Parámetros a ingresar, el tipo de datos y restricciones
            \begin{enumerate}
                \item i. id\_check (parámetro de ruta)
                    \begin{enumerate}
                        \item 1. Number (entero)
                        \item 2. Requerido
                    \end{enumerate}
                \item ii. nombre
                    \begin{enumerate}
                        \item 1. String
                        \item 2. No requerido - Default null
                        \item 3. Tamaño máximo de 50 caracteres
                    \end{enumerate}
                \item iii. fecha\_ter
                    \begin{enumerate}
                        \item 1. Date
                        \item 2. No requerido - Default null
                    \end{enumerate}
            \end{enumerate}
        \item c. Nota: no permite modificar fecha\_ini ni id\_emp
    \end{enumerate}
\end{enumerate}

\subsection{Plantilla de Checklist (/Plantilla)}

\begin{enumerate}
    \item 1. http://”URLPROYECTO”.cl/Plantilla/
    \begin{enumerate}
        \item a. Tipo de petición:
            \begin{itemize}
               \item i. “get”
            \end{itemize}
        \item b. Parámetros a ingresar, el tipo de datos y restricciones
            \begin{enumerate}
                \item i. Sin parámetros de entrada
            \end{enumerate}
    \end{enumerate}

    \item 2. http://”URLPROYECTO”.cl/Plantilla/
    \begin{enumerate}
        \item a. Tipo de petición:
            \begin{itemize}
               \item i. “post”
            \end{itemize}
        \item b. Parámetros a ingresar, el tipo de datos y restricciones
            \begin{enumerate}
                \item i. descr
                    \begin{enumerate}
                        \item 1. String
                        \item 2. No requerido - Default null
                    \end{enumerate}
                \item ii. orden
                    \begin{enumerate}
                        \item 1. Number (entero)
                        \item 2. Requerido
                        \item 3. Debe ser único
                    \end{enumerate}
            \end{enumerate}
    \end{enumerate}

    \item 3. http://”URLPROYECTO”.cl/Plantilla/\{id\_ip\}
    \begin{enumerate}
        \item a. Tipo de petición:
            \begin{itemize}
               \item i. “put”
            \end{itemize}
        \item b. Parámetros a ingresar, el tipo de datos y restricciones
            \begin{enumerate}
                \item i. id\_ip (parámetro de ruta)
                    \begin{enumerate}
                        \item 1. Number (entero)
                        \item 2. Requerido
                    \end{enumerate}
                \item ii. descr
                    \begin{enumerate}
                        \item 1. String
                        \item 2. No requerido - Default null
                    \end{enumerate}
                \item iii. orden
                    \begin{enumerate}
                        \item 1. Number (entero)
                        \item 2. No requerido - Default null
                    \end{enumerate}
            \end{enumerate}
        \item c. Nota: hallazgo técnico — el código referencia una variable antes de definirla, provoca error en tiempo de ejecución al invocarse
    \end{enumerate}
\end{enumerate}

\subsection{Checklist - Plantilla (/check-plantillas)}

\begin{enumerate}
    \item 1. http://”URLPROYECTO”.cl/check-plantillas/
    \begin{enumerate}
        \item a. Tipo de petición:
            \begin{itemize}
               \item i. “get”
            \end{itemize}
        \item b. Parámetros a ingresar, el tipo de datos y restricciones
            \begin{enumerate}
                \item i. Sin parámetros de entrada
            \end{enumerate}
    \end{enumerate}

    \item 2. http://”URLPROYECTO”.cl/check-plantillas/\{id\_check\}/\{id\_ip\}
    \begin{enumerate}
        \item a. Tipo de petición:
            \begin{itemize}
               \item i. “get”
            \end{itemize}
        \item b. Parámetros a ingresar, el tipo de datos y restricciones
            \begin{enumerate}
                \item i. id\_check (parámetro de ruta)
                    \begin{enumerate}
                        \item 1. Number (entero)
                        \item 2. Requerido
                    \end{enumerate}
                \item ii. id\_ip (parámetro de ruta)
                    \begin{enumerate}
                        \item 1. Number (entero)
                        \item 2. Requerido
                    \end{enumerate}
            \end{enumerate}
    \end{enumerate}

    \item 3. http://”URLPROYECTO”.cl/check-plantillas/
    \begin{enumerate}
        \item a. Tipo de petición:
            \begin{itemize}
               \item i. “post”
            \end{itemize}
        \item b. Parámetros a ingresar, el tipo de datos y restricciones
            \begin{enumerate}
                \item i. estado
                    \begin{enumerate}
                        \item 1. String (enum)
                        \item 2. Requerido
                        \item 3. Valores permitidos: "pendiente", "completado", "no\_aplica", "en\_proceso"
                    \end{enumerate}
                \item ii. observacion
                    \begin{enumerate}
                        \item 1. String
                        \item 2. No requerido - Default null
                        \item 3. Tamaño máximo de 255 caracteres
                    \end{enumerate}
                \item iii. id\_check
                    \begin{enumerate}
                        \item 1. Number (entero)
                        \item 2. Requerido
                        \item 3. Debe existir en Checklist
                    \end{enumerate}
                \item iv. id\_ip
                    \begin{enumerate}
                        \item 1. Number (entero)
                        \item 2. Requerido
                        \item 3. Debe existir en Plantilla
                        \item 4. La combinación id\_check + id\_ip no puede estar duplicada
                    \end{enumerate}
            \end{enumerate}
    \end{enumerate}

    \item 4. http://”URLPROYECTO”.cl/check-plantillas/\{id\_check\}/\{id\_ip\}
    \begin{enumerate}
        \item a. Tipo de petición:
            \begin{itemize}
               \item i. “put”
            \end{itemize}
        \item b. Parámetros a ingresar, el tipo de datos y restricciones
            \begin{enumerate}
                \item i. id\_check (parámetro de ruta)
                    \begin{enumerate}
                        \item 1. Number (entero)
                        \item 2. Requerido
                    \end{enumerate}
                \item ii. id\_ip (parámetro de ruta)
                    \begin{enumerate}
                        \item 1. Number (entero)
                        \item 2. Requerido
                    \end{enumerate}
                \item iii. estado
                    \begin{enumerate}
                        \item 1. String (enum)
                        \item 2. No requerido - Default null
                        \item 3. Valores permitidos: "pendiente", "completado", "no\_aplica", "en\_proceso"
                    \end{enumerate}
                \item iv. observacion
                    \begin{enumerate}
                        \item 1. String
                        \item 2. No requerido - Default null
                        \item 3. Tamaño máximo de 255 caracteres
                    \end{enumerate}
            \end{enumerate}
    \end{enumerate}
\end{enumerate}

\subsection{Empresas (/empresas)}

\begin{enumerate}
    \item 1. http://”URLPROYECTO”.cl/empresas/
    \begin{enumerate}
        \item a. Tipo de petición:
            \begin{itemize}
               \item i. “get”
            \end{itemize}
        \item b. Parámetros a ingresar, el tipo de datos y restricciones
            \begin{enumerate}
                \item i. Sin parámetros de entrada
            \end{enumerate}
    \end{enumerate}

    \item 2. http://”URLPROYECTO”.cl/empresas/\{id\_org\}
    \begin{enumerate}
        \item a. Tipo de petición:
            \begin{itemize}
               \item i. “get”
            \end{itemize}
        \item b. Parámetros a ingresar, el tipo de datos y restricciones
            \begin{enumerate}
                \item i. id\_org (parámetro de ruta)
                    \begin{enumerate}
                        \item 1. Number (entero)
                        \item 2. Requerido
                    \end{enumerate}
            \end{enumerate}
    \end{enumerate}

    \item 3. http://”URLPROYECTO”.cl/empresas/
    \begin{enumerate}
        \item a. Tipo de petición:
            \begin{itemize}
               \item i. “post”
            \end{itemize}
        \item b. Parámetros a ingresar, el tipo de datos y restricciones
            \begin{enumerate}
                \item i. nombre
                    \begin{enumerate}
                        \item 1. String
                        \item 2. Requerido
                        \item 3. Tamaño máximo de 100 caracteres
                    \end{enumerate}
            \end{enumerate}
    \end{enumerate}

    \item 4. http://”URLPROYECTO”.cl/empresas/\{id\_org\}
    \begin{enumerate}
        \item a. Tipo de petición:
            \begin{itemize}
               \item i. “put”
            \end{itemize}
        \item b. Parámetros a ingresar, el tipo de datos y restricciones
            \begin{enumerate}
                \item i. id\_org (parámetro de ruta)
                    \begin{enumerate}
                        \item 1. Number (entero)
                        \item 2. Requerido
                    \end{enumerate}
                \item ii. nombre
                    \begin{enumerate}
                        \item 1. String
                        \item 2. No requerido - Default null
                        \item 3. Tamaño máximo de 100 caracteres
                    \end{enumerate}
            \end{enumerate}
    \end{enumerate}

    \item 5. http://”URLPROYECTO”.cl/empresas/\{id\_org\}
    \begin{enumerate}
        \item a. Tipo de petición:
            \begin{itemize}
               \item i. “delete”
            \end{itemize}
        \item b. Parámetros a ingresar, el tipo de datos y restricciones
            \begin{enumerate}
                \item i. id\_org (parámetro de ruta)
                    \begin{enumerate}
                        \item 1. Number (entero)
                        \item 2. Requerido
                        \item 3. Bloqueado (400) si tiene Departamentos asociados
                    \end{enumerate}
            \end{enumerate}
        \item c. Nota: borrado físico
    \end{enumerate}
\end{enumerate}

\subsection{Cargos (/cargos)}

\begin{enumerate}
    \item 1. http://”URLPROYECTO”.cl/cargos/
    \begin{enumerate}
        \item a. Tipo de petición:
            \begin{itemize}
               \item i. “get”
            \end{itemize}
        \item b. Parámetros a ingresar, el tipo de datos y restricciones
            \begin{enumerate}
                \item i. Sin parámetros de entrada
            \end{enumerate}
    \end{enumerate}

    \item 2. http://”URLPROYECTO”.cl/cargos/\{id\_cargo\}
    \begin{enumerate}
        \item a. Tipo de petición:
            \begin{itemize}
               \item i. “get”
            \end{itemize}
        \item b. Parámetros a ingresar, el tipo de datos y restricciones
            \begin{enumerate}
                \item i. id\_cargo (parámetro de ruta)
                    \begin{enumerate}
                        \item 1. Number (entero)
                        \item 2. Requerido
                    \end{enumerate}
            \end{enumerate}
    \end{enumerate}

    \item 3. http://”URLPROYECTO”.cl/cargos/
    \begin{enumerate}
        \item a. Tipo de petición:
            \begin{itemize}
               \item i. “post”
            \end{itemize}
        \item b. Parámetros a ingresar, el tipo de datos y restricciones
            \begin{enumerate}
                \item i. nombre
                    \begin{enumerate}
                        \item 1. String
                        \item 2. Requerido
                        \item 3. Tamaño máximo de 100 caracteres
                    \end{enumerate}
            \end{enumerate}
    \end{enumerate}

    \item 4. http://”URLPROYECTO”.cl/cargos/\{id\_cargo\}
    \begin{enumerate}
        \item a. Tipo de petición:
            \begin{itemize}
               \item i. “put”
            \end{itemize}
        \item b. Parámetros a ingresar, el tipo de datos y restricciones
            \begin{enumerate}
                \item i. id\_cargo (parámetro de ruta)
                    \begin{enumerate}
                        \item 1. Number (entero)
                        \item 2. Requerido
                    \end{enumerate}
                \item ii. nombre
                    \begin{enumerate}
                        \item 1. String
                        \item 2. No requerido - Default null
                        \item 3. Tamaño máximo de 100 caracteres
                    \end{enumerate}
            \end{enumerate}
    \end{enumerate}

    \item 5. http://”URLPROYECTO”.cl/cargos/\{id\_cargo\}
    \begin{enumerate}
        \item a. Tipo de petición:
            \begin{itemize}
               \item i. “delete”
            \end{itemize}
        \item b. Parámetros a ingresar, el tipo de datos y restricciones
            \begin{enumerate}
                \item i. id\_cargo (parámetro de ruta)
                    \begin{enumerate}
                        \item 1. Number (entero)
                        \item 2. Requerido
                        \item 3. Bloqueado (400) si tiene Empleados asociados
                    \end{enumerate}
            \end{enumerate}
        \item c. Nota: borrado físico
    \end{enumerate}
\end{enumerate}

\subsection{Plantas (/plantas)}

\begin{enumerate}
    \item 1. http://”URLPROYECTO”.cl/plantas/
    \begin{enumerate}
        \item a. Tipo de petición:
            \begin{itemize}
               \item i. “get”
            \end{itemize}
        \item b. Parámetros a ingresar, el tipo de datos y restricciones
            \begin{enumerate}
                \item i. Sin parámetros de entrada
            \end{enumerate}
    \end{enumerate}

    \item 2. http://”URLPROYECTO”.cl/plantas/\{id\_planta\}
    \begin{enumerate}
        \item a. Tipo de petición:
            \begin{itemize}
               \item i. “get”
            \end{itemize}
        \item b. Parámetros a ingresar, el tipo de datos y restricciones
            \begin{enumerate}
                \item i. id\_planta (parámetro de ruta)
                    \begin{enumerate}
                        \item 1. Number (entero)
                        \item 2. Requerido
                    \end{enumerate}
            \end{enumerate}
    \end{enumerate}

    \item 3. http://”URLPROYECTO”.cl/plantas/
    \begin{enumerate}
        \item a. Tipo de petición:
            \begin{itemize}
               \item i. “post”
            \end{itemize}
        \item b. Parámetros a ingresar, el tipo de datos y restricciones
            \begin{enumerate}
                \item i. nombre
                    \begin{enumerate}
                        \item 1. String
                        \item 2. Requerido
                        \item 3. Tamaño máximo de 100 caracteres
                    \end{enumerate}
            \end{enumerate}
    \end{enumerate}

    \item 4. http://”URLPROYECTO”.cl/plantas/\{id\_planta\}
    \begin{enumerate}
        \item a. Tipo de petición:
            \begin{itemize}
               \item i. “put”
            \end{itemize}
        \item b. Parámetros a ingresar, el tipo de datos y restricciones
            \begin{enumerate}
                \item i. id\_planta (parámetro de ruta)
                    \begin{enumerate}
                        \item 1. Number (entero)
                        \item 2. Requerido
                    \end{enumerate}
                \item ii. nombre
                    \begin{enumerate}
                        \item 1. String
                        \item 2. No requerido - Default null
                        \item 3. Tamaño máximo de 100 caracteres
                    \end{enumerate}
            \end{enumerate}
        \item c. Nota: hallazgo técnico — el código referencia variables antes de definirlas, provoca error en tiempo de ejecución al invocarse
    \end{enumerate}

    \item 5. http://”URLPROYECTO”.cl/plantas/\{id\_planta\}
    \begin{enumerate}
        \item a. Tipo de petición:
            \begin{itemize}
               \item i. “delete”
            \end{itemize}
        \item b. Parámetros a ingresar, el tipo de datos y restricciones
            \begin{enumerate}
                \item i. id\_planta (parámetro de ruta)
                    \begin{enumerate}
                        \item 1. Number (entero)
                        \item 2. Requerido
                        \item 3. Bloqueado (400) si tiene Empleados asociados
                    \end{enumerate}
            \end{enumerate}
        \item c. Nota: borrado físico
    \end{enumerate}
\end{enumerate}

\subsection{Tipos de Equipo (/tipos-equipo)}

\begin{enumerate}
    \item 1. http://”URLPROYECTO”.cl/tipos-equipo/
    \begin{enumerate}
        \item a. Tipo de petición:
            \begin{itemize}
               \item i. “get”
            \end{itemize}
        \item b. Parámetros a ingresar, el tipo de datos y restricciones
            \begin{enumerate}
                \item i. Sin parámetros de entrada
            \end{enumerate}
    \end{enumerate}

    \item 2. http://”URLPROYECTO”.cl/tipos-equipo/\{id\_tipoeq\}
    \begin{enumerate}
        \item a. Tipo de petición:
            \begin{itemize}
               \item i. “get”
            \end{itemize}
        \item b. Parámetros a ingresar, el tipo de datos y restricciones
            \begin{enumerate}
                \item i. id\_tipoeq (parámetro de ruta)
                    \begin{enumerate}
                        \item 1. Number (entero)
                        \item 2. Requerido
                    \end{enumerate}
            \end{enumerate}
    \end{enumerate}

    \item 3. http://”URLPROYECTO”.cl/tipos-equipo/
    \begin{enumerate}
        \item a. Tipo de petición:
            \begin{itemize}
               \item i. “post”
            \end{itemize}
        \item b. Parámetros a ingresar, el tipo de datos y restricciones
            \begin{enumerate}
                \item i. nombre
                    \begin{enumerate}
                        \item 1. String
                        \item 2. Requerido
                        \item 3. Tamaño máximo de 100 caracteres
                    \end{enumerate}
            \end{enumerate}
    \end{enumerate}

    \item 4. http://”URLPROYECTO”.cl/tipos-equipo/\{id\_tipoeq\}
    \begin{enumerate}
        \item a. Tipo de petición:
            \begin{itemize}
               \item i. “put”
            \end{itemize}
        \item b. Parámetros a ingresar, el tipo de datos y restricciones
            \begin{enumerate}
                \item i. id\_tipoeq (parámetro de ruta)
                    \begin{enumerate}
                        \item 1. Number (entero)
                        \item 2. Requerido
                    \end{enumerate}
                \item ii. nombre
                    \begin{enumerate}
                        \item 1. String
                        \item 2. No requerido - Default null
                        \item 3. Tamaño máximo de 100 caracteres
                    \end{enumerate}
            \end{enumerate}
    \end{enumerate}

    \item 5. http://”URLPROYECTO”.cl/tipos-equipo/\{id\_tipoeq\}
    \begin{enumerate}
        \item a. Tipo de petición:
            \begin{itemize}
               \item i. “delete”
            \end{itemize}
        \item b. Parámetros a ingresar, el tipo de datos y restricciones
            \begin{enumerate}
                \item i. id\_tipoeq (parámetro de ruta)
                    \begin{enumerate}
                        \item 1. Number (entero)
                        \item 2. Requerido
                        \item 3. Bloqueado (400) si tiene Equipos de Cómputo asociados
                    \end{enumerate}
            \end{enumerate}
        \item c. Nota: borrado físico
    \end{enumerate}
\end{enumerate}

\subsection{Sistemas Operativos (/sistemas-operativos) - Familia}

\begin{enumerate}
    \item 1. http://”URLPROYECTO”.cl/sistemas-operativos/familias
    \begin{enumerate}
        \item a. Tipo de petición:
            \begin{itemize}
               \item i. “get”
            \end{itemize}
        \item b. Parámetros a ingresar, el tipo de datos y restricciones
            \begin{enumerate}
                \item i. Sin parámetros de entrada
            \end{enumerate}
    \end{enumerate}

    \item 2. http://”URLPROYECTO”.cl/sistemas-operativos/familias
    \begin{enumerate}
        \item a. Tipo de petición:
            \begin{itemize}
               \item i. “post”
            \end{itemize}
        \item b. Parámetros a ingresar, el tipo de datos y restricciones
            \begin{enumerate}
                \item i. nombre
                    \begin{enumerate}
                        \item 1. String
                        \item 2. Requerido
                        \item 3. Tamaño máximo de 100 caracteres
                        \item 4. Debe ser único
                    \end{enumerate}
            \end{enumerate}
    \end{enumerate}

    \item 3. http://”URLPROYECTO”.cl/sistemas-operativos/familias/\{id\_fam\_so\}
    \begin{enumerate}
        \item a. Tipo de petición:
            \begin{itemize}
               \item i. “put”
            \end{itemize}
        \item b. Parámetros a ingresar, el tipo de datos y restricciones
            \begin{enumerate}
                \item i. id\_fam\_so (parámetro de ruta)
                    \begin{enumerate}
                        \item 1. Number (entero)
                        \item 2. Requerido
                    \end{enumerate}
                \item ii. nombre
                    \begin{enumerate}
                        \item 1. String
                        \item 2. No requerido - Default null
                        \item 3. Tamaño máximo de 100 caracteres
                    \end{enumerate}
            \end{enumerate}
    \end{enumerate}

    \item 4. http://”URLPROYECTO”.cl/sistemas-operativos/familias/\{id\_fam\_so\}
    \begin{enumerate}
        \item a. Tipo de petición:
            \begin{itemize}
               \item i. “delete”
            \end{itemize}
        \item b. Parámetros a ingresar, el tipo de datos y restricciones
            \begin{enumerate}
                \item i. id\_fam\_so (parámetro de ruta)
                    \begin{enumerate}
                        \item 1. Number (entero)
                        \item 2. Requerido
                        \item 3. Bloqueado (400) si tiene Versiones de SO asociadas
                    \end{enumerate}
            \end{enumerate}
        \item c. Nota: borrado físico
    \end{enumerate}
\end{enumerate}

\subsection{Sistemas Operativos (/sistemas-operativos) - Versión}

\begin{enumerate}
    \item 1. http://”URLPROYECTO”.cl/sistemas-operativos/versiones
    \begin{enumerate}
        \item a. Tipo de petición:
            \begin{itemize}
               \item i. “get”
            \end{itemize}
        \item b. Parámetros a ingresar, el tipo de datos y restricciones
            \begin{enumerate}
                \item i. Sin parámetros de entrada
            \end{enumerate}
    \end{enumerate}

    \item 2. http://”URLPROYECTO”.cl/sistemas-operativos/versiones/por-familia/\{id\_fam\_so\}
    \begin{enumerate}
        \item a. Tipo de petición:
            \begin{itemize}
               \item i. “get”
            \end{itemize}
        \item b. Parámetros a ingresar, el tipo de datos y restricciones
            \begin{enumerate}
                \item i. id\_fam\_so (parámetro de ruta)
                    \begin{enumerate}
                        \item 1. Number (entero)
                        \item 2. Requerido
                    \end{enumerate}
            \end{enumerate}
    \end{enumerate}

    \item 3. http://”URLPROYECTO”.cl/sistemas-operativos/versiones
    \begin{enumerate}
        \item a. Tipo de petición:
            \begin{itemize}
               \item i. “post”
            \end{itemize}
        \item b. Parámetros a ingresar, el tipo de datos y restricciones
            \begin{enumerate}
                \item i. nombre
                    \begin{enumerate}
                        \item 1. String
                        \item 2. Requerido
                        \item 3. Tamaño máximo de 100 caracteres
                    \end{enumerate}
                \item ii. id\_fam\_so
                    \begin{enumerate}
                        \item 1. Number (entero)
                        \item 2. Requerido
                        \item 3. Debe existir en Familia de SO
                    \end{enumerate}
            \end{enumerate}
    \end{enumerate}

    \item 4. http://”URLPROYECTO”.cl/sistemas-operativos/versiones/\{id\_ver\_so\}
    \begin{enumerate}
        \item a. Tipo de petición:
            \begin{itemize}
               \item i. “put”
            \end{itemize}
        \item b. Parámetros a ingresar, el tipo de datos y restricciones
            \begin{enumerate}
                \item i. id\_ver\_so (parámetro de ruta)
                    \begin{enumerate}
                        \item 1. Number (entero)
                        \item 2. Requerido
                    \end{enumerate}
                \item ii. nombre
                    \begin{enumerate}
                        \item 1. String
                        \item 2. No requerido - Default null
                        \item 3. Tamaño máximo de 100 caracteres
                    \end{enumerate}
                \item iii. id\_fam\_so
                    \begin{enumerate}
                        \item 1. Number (entero)
                        \item 2. No requerido - Default null
                    \end{enumerate}
            \end{enumerate}
    \end{enumerate}

    \item 5. http://”URLPROYECTO”.cl/sistemas-operativos/versiones/\{id\_ver\_so\}
    \begin{enumerate}
        \item a. Tipo de petición:
            \begin{itemize}
               \item i. “delete”
            \end{itemize}
        \item b. Parámetros a ingresar, el tipo de datos y restricciones
            \begin{enumerate}
                \item i. id\_ver\_so (parámetro de ruta)
                    \begin{enumerate}
                        \item 1. Number (entero)
                        \item 2. Requerido
                        \item 3. Bloqueado (400) si tiene Ediciones de SO asociadas
                    \end{enumerate}
            \end{enumerate}
        \item c. Nota: borrado físico
    \end{enumerate}
\end{enumerate}

\subsection{Sistemas Operativos (/sistemas-operativos) - Edición}

\begin{enumerate}
    \item 1. http://”URLPROYECTO”.cl/sistemas-operativos/ediciones
    \begin{enumerate}
        \item a. Tipo de petición:
            \begin{itemize}
               \item i. “get”
            \end{itemize}
        \item b. Parámetros a ingresar, el tipo de datos y restricciones
            \begin{enumerate}
                \item i. Sin parámetros de entrada
            \end{enumerate}
    \end{enumerate}

    \item 2. http://”URLPROYECTO”.cl/sistemas-operativos/ediciones/por-version/\{id\_ver\_so\}
    \begin{enumerate}
        \item a. Tipo de petición:
            \begin{itemize}
               \item i. “get”
            \end{itemize}
        \item b. Parámetros a ingresar, el tipo de datos y restricciones
            \begin{enumerate}
                \item i. id\_ver\_so (parámetro de ruta)
                    \begin{enumerate}
                        \item 1. Number (entero)
                        \item 2. Requerido
                    \end{enumerate}
            \end{enumerate}
    \end{enumerate}

    \item 3. http://”URLPROYECTO”.cl/sistemas-operativos/ediciones
    \begin{enumerate}
        \item a. Tipo de petición:
            \begin{itemize}
               \item i. “post”
            \end{itemize}
        \item b. Parámetros a ingresar, el tipo de datos y restricciones
            \begin{enumerate}
                \item i. nombre
                    \begin{enumerate}
                        \item 1. String
                        \item 2. Requerido
                        \item 3. Tamaño máximo de 100 caracteres
                    \end{enumerate}
                \item ii. id\_ver\_so
                    \begin{enumerate}
                        \item 1. Number (entero)
                        \item 2. Requerido
                        \item 3. Debe existir en Versión de SO
                    \end{enumerate}
            \end{enumerate}
    \end{enumerate}

    \item 4. http://”URLPROYECTO”.cl/sistemas-operativos/ediciones/\{id\_edicion\_so\}
    \begin{enumerate}
        \item a. Tipo de petición:
            \begin{itemize}
               \item i. “put”
            \end{itemize}
        \item b. Parámetros a ingresar, el tipo de datos y restricciones
            \begin{enumerate}
                \item i. id\_edicion\_so (parámetro de ruta)
                    \begin{enumerate}
                        \item 1. Number (entero)
                        \item 2. Requerido
                    \end{enumerate}
                \item ii. nombre
                    \begin{enumerate}
                        \item 1. String
                        \item 2. No requerido - Default null
                        \item 3. Tamaño máximo de 100 caracteres
                    \end{enumerate}
                \item iii. id\_ver\_so
                    \begin{enumerate}
                        \item 1. Number (entero)
                        \item 2. No requerido - Default null
                    \end{enumerate}
            \end{enumerate}
    \end{enumerate}

    \item 5. http://”URLPROYECTO”.cl/sistemas-operativos/ediciones/\{id\_edicion\_so\}
    \begin{enumerate}
        \item a. Tipo de petición:
            \begin{itemize}
               \item i. “delete”
            \end{itemize}
        \item b. Parámetros a ingresar, el tipo de datos y restricciones
            \begin{enumerate}
                \item i. id\_edicion\_so (parámetro de ruta)
                    \begin{enumerate}
                        \item 1. Number (entero)
                        \item 2. Requerido
                        \item 3. Bloqueado (400) si tiene Equipos de Cómputo asociados
                    \end{enumerate}
            \end{enumerate}
        \item c. Nota: borrado físico
    \end{enumerate}
\end{enumerate}

\subsection{Tipos de Periférico (/tipos-periferico)}

\begin{enumerate}
    \item 1. http://”URLPROYECTO”.cl/tipos-periferico/
    \begin{enumerate}
        \item a. Tipo de petición:
            \begin{itemize}
               \item i. “get”
            \end{itemize}
        \item b. Parámetros a ingresar, el tipo de datos y restricciones
            \begin{enumerate}
                \item i. Sin parámetros de entrada
            \end{enumerate}
    \end{enumerate}

    \item 2. http://”URLPROYECTO”.cl/tipos-periferico/\{id\_tipoper\}
    \begin{enumerate}
        \item a. Tipo de petición:
            \begin{itemize}
               \item i. “get”
            \end{itemize}
        \item b. Parámetros a ingresar, el tipo de datos y restricciones
            \begin{enumerate}
                \item i. id\_tipoper (parámetro de ruta)
                    \begin{enumerate}
                        \item 1. Number (entero)
                        \item 2. Requerido
                    \end{enumerate}
            \end{enumerate}
    \end{enumerate}

    \item 3. http://”URLPROYECTO”.cl/tipos-periferico/
    \begin{enumerate}
        \item a. Tipo de petición:
            \begin{itemize}
               \item i. “post”
            \end{itemize}
        \item b. Parámetros a ingresar, el tipo de datos y restricciones
            \begin{enumerate}
                \item i. nombre
                    \begin{enumerate}
                        \item 1. String
                        \item 2. Requerido
                        \item 3. Tamaño máximo de 100 caracteres
                    \end{enumerate}
            \end{enumerate}
    \end{enumerate}

    \item 4. http://”URLPROYECTO”.cl/tipos-periferico/\{id\_tipoper\}
    \begin{enumerate}
        \item a. Tipo de petición:
            \begin{itemize}
               \item i. “put”
            \end{itemize}
        \item b. Parámetros a ingresar, el tipo de datos y restricciones
            \begin{enumerate}
                \item i. id\_tipoper (parámetro de ruta)
                    \begin{enumerate}
                        \item 1. Number (entero)
                        \item 2. Requerido
                    \end{enumerate}
                \item ii. nombre
                    \begin{enumerate}
                        \item 1. String
                        \item 2. No requerido - Default null
                        \item 3. Tamaño máximo de 100 caracteres
                    \end{enumerate}
            \end{enumerate}
    \end{enumerate}

    \item 5. http://”URLPROYECTO”.cl/tipos-periferico/\{id\_tipoper\}
    \begin{enumerate}
        \item a. Tipo de petición:
            \begin{itemize}
               \item i. “delete”
            \end{itemize}
        \item b. Parámetros a ingresar, el tipo de datos y restricciones
            \begin{enumerate}
                \item i. id\_tipoper (parámetro de ruta)
                    \begin{enumerate}
                        \item 1. Number (entero)
                        \item 2. Requerido
                        \item 3. Bloqueado (400) si tiene Periféricos asociados
                    \end{enumerate}
            \end{enumerate}
        \item c. Nota: borrado físico
    \end{enumerate}
\end{enumerate}

\subsection{Proveedores de Licencia (/proveedores-licencia)}

\begin{enumerate}
    \item 1. http://”URLPROYECTO”.cl/proveedores-licencia/
    \begin{enumerate}
        \item a. Tipo de petición:
            \begin{itemize}
               \item i. “get”
            \end{itemize}
        \item b. Parámetros a ingresar, el tipo de datos y restricciones
            \begin{enumerate}
                \item i. Sin parámetros de entrada
            \end{enumerate}
    \end{enumerate}

    \item 2. http://”URLPROYECTO”.cl/proveedores-licencia/\{id\_prov\}
    \begin{enumerate}
        \item a. Tipo de petición:
            \begin{itemize}
               \item i. “get”
            \end{itemize}
        \item b. Parámetros a ingresar, el tipo de datos y restricciones
            \begin{enumerate}
                \item i. id\_prov (parámetro de ruta)
                    \begin{enumerate}
                        \item 1. Number (entero)
                        \item 2. Requerido
                    \end{enumerate}
            \end{enumerate}
    \end{enumerate}

    \item 3. http://”URLPROYECTO”.cl/proveedores-licencia/
    \begin{enumerate}
        \item a. Tipo de petición:
            \begin{itemize}
               \item i. “post”
            \end{itemize}
        \item b. Parámetros a ingresar, el tipo de datos y restricciones
            \begin{enumerate}
                \item i. nombre
                    \begin{enumerate}
                        \item 1. String
                        \item 2. Requerido
                        \item 3. Tamaño máximo de 100 caracteres
                    \end{enumerate}
            \end{enumerate}
    \end{enumerate}

    \item 4. http://”URLPROYECTO”.cl/proveedores-licencia/\{id\_prov\}
    \begin{enumerate}
        \item a. Tipo de petición:
            \begin{itemize}
               \item i. “put”
            \end{itemize}
        \item b. Parámetros a ingresar, el tipo de datos y restricciones
            \begin{enumerate}
                \item i. id\_prov (parámetro de ruta)
                    \begin{enumerate}
                        \item 1. Number (entero)
                        \item 2. Requerido
                    \end{enumerate}
                \item ii. nombre
                    \begin{enumerate}
                        \item 1. String
                        \item 2. No requerido - Default null
                        \item 3. Tamaño máximo de 100 caracteres
                    \end{enumerate}
            \end{enumerate}
    \end{enumerate}

    \item 5. http://”URLPROYECTO”.cl/proveedores-licencia/\{id\_prov\}
    \begin{enumerate}
        \item a. Tipo de petición:
            \begin{itemize}
               \item i. “delete”
            \end{itemize}
        \item b. Parámetros a ingresar, el tipo de datos y restricciones
            \begin{enumerate}
                \item i. id\_prov (parámetro de ruta)
                    \begin{enumerate}
                        \item 1. Number (entero)
                        \item 2. Requerido
                        \item 3. Bloqueado (400) si tiene Versiones de Licencia asociadas
                    \end{enumerate}
            \end{enumerate}
        \item c. Nota: borrado físico
    \end{enumerate}
\end{enumerate}

\subsection{Proveedores de Factura (/proveedores-factura)}

\begin{enumerate}
    \item 1. http://”URLPROYECTO”.cl/proveedores-factura/
    \begin{enumerate}
        \item a. Tipo de petición:
            \begin{itemize}
               \item i. “get”
            \end{itemize}
        \item b. Parámetros a ingresar, el tipo de datos y restricciones
            \begin{enumerate}
                \item i. Sin parámetros de entrada
            \end{enumerate}
    \end{enumerate}

    \item 2. http://”URLPROYECTO”.cl/proveedores-factura/\{id\_prov\_fact\}
    \begin{enumerate}
        \item a. Tipo de petición:
            \begin{itemize}
               \item i. “get”
            \end{itemize}
        \item b. Parámetros a ingresar, el tipo de datos y restricciones
            \begin{enumerate}
                \item i. id\_prov\_fact (parámetro de ruta)
                    \begin{enumerate}
                        \item 1. Number (entero)
                        \item 2. Requerido
                    \end{enumerate}
            \end{enumerate}
    \end{enumerate}

    \item 3. http://”URLPROYECTO”.cl/proveedores-factura/
    \begin{enumerate}
        \item a. Tipo de petición:
            \begin{itemize}
               \item i. “post”
            \end{itemize}
        \item b. Parámetros a ingresar, el tipo de datos y restricciones
            \begin{enumerate}
                \item i. nombre
                    \begin{enumerate}
                        \item 1. String
                        \item 2. Requerido
                        \item 3. Tamaño máximo de 100 caracteres
                    \end{enumerate}
            \end{enumerate}
    \end{enumerate}

    \item 4. http://”URLPROYECTO”.cl/proveedores-factura/\{id\_prov\_fact\}
    \begin{enumerate}
        \item a. Tipo de petición:
            \begin{itemize}
               \item i. “put”
            \end{itemize}
        \item b. Parámetros a ingresar, el tipo de datos y restricciones
            \begin{enumerate}
                \item i. id\_prov\_fact (parámetro de ruta)
                    \begin{enumerate}
                        \item 1. Number (entero)
                        \item 2. Requerido
                    \end{enumerate}
                \item ii. nombre
                    \begin{enumerate}
                        \item 1. String
                        \item 2. No requerido - Default null
                        \item 3. Tamaño máximo de 100 caracteres
                    \end{enumerate}
            \end{enumerate}
    \end{enumerate}

    \item 5. http://”URLPROYECTO”.cl/proveedores-factura/\{id\_prov\_fact\}
    \begin{enumerate}
        \item a. Tipo de petición:
            \begin{itemize}
               \item i. “delete”
            \end{itemize}
        \item b. Parámetros a ingresar, el tipo de datos y restricciones
            \begin{enumerate}
                \item i. id\_prov\_fact (parámetro de ruta)
                    \begin{enumerate}
                        \item 1. Number (entero)
                        \item 2. Requerido
                        \item 3. Bloqueado (400) si tiene Facturas asociadas
                    \end{enumerate}
            \end{enumerate}
        \item c. Nota: borrado físico
    \end{enumerate}
\end{enumerate}

\subsection{Versiones de Licencia (/versiones-licencia)}

\begin{enumerate}
    \item 1. http://”URLPROYECTO”.cl/versiones-licencia/
    \begin{enumerate}
        \item a. Tipo de petición:
            \begin{itemize}
               \item i. “get”
            \end{itemize}
        \item b. Parámetros a ingresar, el tipo de datos y restricciones
            \begin{enumerate}
                \item i. id\_prov (parámetro de query)
                    \begin{enumerate}
                        \item 1. Number (entero)
                        \item 2. No requerido - Default null
                    \end{enumerate}
            \end{enumerate}
    \end{enumerate}

    \item 2. http://”URLPROYECTO”.cl/versiones-licencia/\{id\_ver\}
    \begin{enumerate}
        \item a. Tipo de petición:
            \begin{itemize}
               \item i. “get”
            \end{itemize}
        \item b. Parámetros a ingresar, el tipo de datos y restricciones
            \begin{enumerate}
                \item i. id\_ver (parámetro de ruta)
                    \begin{enumerate}
                        \item 1. Number (entero)
                        \item 2. Requerido
                    \end{enumerate}
            \end{enumerate}
    \end{enumerate}

    \item 3. http://”URLPROYECTO”.cl/versiones-licencia/
    \begin{enumerate}
        \item a. Tipo de petición:
            \begin{itemize}
               \item i. “post”
            \end{itemize}
        \item b. Parámetros a ingresar, el tipo de datos y restricciones
            \begin{enumerate}
                \item i. nombre
                    \begin{enumerate}
                        \item 1. String
                        \item 2. Requerido
                        \item 3. Tamaño máximo de 100 caracteres
                    \end{enumerate}
                \item ii. id\_prov
                    \begin{enumerate}
                        \item 1. Number (entero)
                        \item 2. Requerido
                        \item 3. Debe existir en Proveedor de Licencia
                    \end{enumerate}
            \end{enumerate}
    \end{enumerate}

    \item 4. http://”URLPROYECTO”.cl/versiones-licencia/\{id\_ver\}
    \begin{enumerate}
        \item a. Tipo de petición:
            \begin{itemize}
               \item i. “put”
            \end{itemize}
        \item b. Parámetros a ingresar, el tipo de datos y restricciones
            \begin{enumerate}
                \item i. id\_ver (parámetro de ruta)
                    \begin{enumerate}
                        \item 1. Number (entero)
                        \item 2. Requerido
                    \end{enumerate}
                \item ii. nombre
                    \begin{enumerate}
                        \item 1. String
                        \item 2. No requerido - Default null
                        \item 3. Tamaño máximo de 100 caracteres
                    \end{enumerate}
                \item iii. id\_prov
                    \begin{enumerate}
                        \item 1. Number (entero)
                        \item 2. No requerido - Default null
                    \end{enumerate}
            \end{enumerate}
    \end{enumerate}

    \item 5. http://”URLPROYECTO”.cl/versiones-licencia/\{id\_ver\}
    \begin{enumerate}
        \item a. Tipo de petición:
            \begin{itemize}
               \item i. “delete”
            \end{itemize}
        \item b. Parámetros a ingresar, el tipo de datos y restricciones
            \begin{enumerate}
                \item i. id\_ver (parámetro de ruta)
                    \begin{enumerate}
                        \item 1. Number (entero)
                        \item 2. Requerido
                        \item 3. Bloqueado (400) si tiene Licencias asociadas
                    \end{enumerate}
            \end{enumerate}
        \item c. Nota: borrado físico
    \end{enumerate}
\end{enumerate}

\subsection{Sub-áreas (/sub-areas)}

\begin{enumerate}
    \item 1. http://”URLPROYECTO”.cl/sub-areas/
    \begin{enumerate}
        \item a. Tipo de petición:
            \begin{itemize}
               \item i. “get”
            \end{itemize}
        \item b. Parámetros a ingresar, el tipo de datos y restricciones
            \begin{enumerate}
                \item i. id\_dpto (parámetro de query)
                    \begin{enumerate}
                        \item 1. Number (entero)
                        \item 2. No requerido - Default null
                    \end{enumerate}
            \end{enumerate}
    \end{enumerate}

    \item 2. http://”URLPROYECTO”.cl/sub-areas/\{id\_sub\}
    \begin{enumerate}
        \item a. Tipo de petición:
            \begin{itemize}
               \item i. “get”
            \end{itemize}
        \item b. Parámetros a ingresar, el tipo de datos y restricciones
            \begin{enumerate}
                \item i. id\_sub (parámetro de ruta)
                    \begin{enumerate}
                        \item 1. Number (entero)
                        \item 2. Requerido
                    \end{enumerate}
            \end{enumerate}
    \end{enumerate}

    \item 3. http://”URLPROYECTO”.cl/sub-areas/
    \begin{enumerate}
        \item a. Tipo de petición:
            \begin{itemize}
               \item i. “post”
            \end{itemize}
        \item b. Parámetros a ingresar, el tipo de datos y restricciones
            \begin{enumerate}
                \item i. nombre
                    \begin{enumerate}
                        \item 1. String
                        \item 2. Requerido
                        \item 3. Tamaño máximo de 100 caracteres
                    \end{enumerate}
                \item ii. id\_dpto
                    \begin{enumerate}
                        \item 1. Number (entero)
                        \item 2. Requerido
                        \item 3. Debe existir en Departamento
                    \end{enumerate}
            \end{enumerate}
        \item c. Nota: este endpoint no exige autenticación (posible pendiente de seguridad)
    \end{enumerate}

    \item 4. http://”URLPROYECTO”.cl/sub-areas/\{id\_sub\}
    \begin{enumerate}
        \item a. Tipo de petición:
            \begin{itemize}
               \item i. “put”
            \end{itemize}
        \item b. Parámetros a ingresar, el tipo de datos y restricciones
            \begin{enumerate}
                \item i. id\_sub (parámetro de ruta)
                    \begin{enumerate}
                        \item 1. Number (entero)
                        \item 2. Requerido
                    \end{enumerate}
                \item ii. nombre
                    \begin{enumerate}
                        \item 1. String
                        \item 2. No requerido - Default null
                        \item 3. Tamaño máximo de 100 caracteres
                    \end{enumerate}
                \item iii. id\_dpto
                    \begin{enumerate}
                        \item 1. Number (entero)
                        \item 2. No requerido - Default null
                    \end{enumerate}
            \end{enumerate}
    \end{enumerate}

    \item 5. http://”URLPROYECTO”.cl/sub-areas/\{id\_sub\}
    \begin{enumerate}
        \item a. Tipo de petición:
            \begin{itemize}
               \item i. “delete”
            \end{itemize}
        \item b. Parámetros a ingresar, el tipo de datos y restricciones
            \begin{enumerate}
                \item i. id\_sub (parámetro de ruta)
                    \begin{enumerate}
                        \item 1. Number (entero)
                        \item 2. Requerido
                        \item 3. Bloqueado (400) si tiene Empleados asociados
                    \end{enumerate}
            \end{enumerate}
        \item c. Nota: borrado físico
    \end{enumerate}
\end{enumerate}

\subsection{Disponibilidades (/disponibilidades)}

\begin{enumerate}
    \item 1. http://”URLPROYECTO”.cl/disponibilidades/
    \begin{enumerate}
        \item a. Tipo de petición:
            \begin{itemize}
               \item i. “get”
            \end{itemize}
        \item b. Parámetros a ingresar, el tipo de datos y restricciones
            \begin{enumerate}
                \item i. Sin parámetros de entrada
            \end{enumerate}
    \end{enumerate}

    \item 2. http://”URLPROYECTO”.cl/disponibilidades/
    \begin{enumerate}
        \item a. Tipo de petición:
            \begin{itemize}
               \item i. “post”
            \end{itemize}
        \item b. Parámetros a ingresar, el tipo de datos y restricciones
            \begin{enumerate}
                \item i. nombre
                    \begin{enumerate}
                        \item 1. String
                        \item 2. Requerido
                        \item 3. Debe ser único
                    \end{enumerate}
                \item ii. descr
                    \begin{enumerate}
                        \item 1. String
                        \item 2. Requerido
                    \end{enumerate}
                \item iii. color
                    \begin{enumerate}
                        \item 1. String
                        \item 2. No requerido - Default null
                        \item 3. Valores esperados: "success", "warning", "error", "info", "primary", "default" (no aplicado como enum estricto)
                        \item 4. Tamaño máximo de 20 caracteres
                    \end{enumerate}
            \end{enumerate}
        \item c. Nota: este endpoint requiere solo autenticación (no exige rol editor)
    \end{enumerate}

    \item 3. http://”URLPROYECTO”.cl/disponibilidades/\{id\_dispo\}
    \begin{enumerate}
        \item a. Tipo de petición:
            \begin{itemize}
               \item i. “put”
            \end{itemize}
        \item b. Parámetros a ingresar, el tipo de datos y restricciones
            \begin{enumerate}
                \item i. id\_dispo (parámetro de ruta)
                    \begin{enumerate}
                        \item 1. Number (entero)
                        \item 2. Requerido
                    \end{enumerate}
                \item ii. nombre, descr, color
                    \begin{enumerate}
                        \item 1. Mismos tipos y restricciones que en la creación
                        \item 2. No requeridos - Default null
                    \end{enumerate}
            \end{enumerate}
    \end{enumerate}

    \item 4. http://”URLPROYECTO”.cl/disponibilidades/\{id\_dispo\}
    \begin{enumerate}
        \item a. Tipo de petición:
            \begin{itemize}
               \item i. “delete”
            \end{itemize}
        \item b. Parámetros a ingresar, el tipo de datos y restricciones
            \begin{enumerate}
                \item i. id\_dispo (parámetro de ruta)
                    \begin{enumerate}
                        \item 1. Number (entero)
                        \item 2. Requerido
                    \end{enumerate}
            \end{enumerate}
        \item c. Nota: borrado físico
    \end{enumerate}

    \item 5. http://”URLPROYECTO”.cl/disponibilidades/por-estado/\{id\_tipoe\}
    \begin{enumerate}
        \item a. Tipo de petición:
            \begin{itemize}
               \item i. “get”
            \end{itemize}
        \item b. Parámetros a ingresar, el tipo de datos y restricciones
            \begin{enumerate}
                \item i. id\_tipoe (parámetro de ruta)
                    \begin{enumerate}
                        \item 1. Number (entero)
                        \item 2. Requerido
                    \end{enumerate}
            \end{enumerate}
    \end{enumerate}

    \item 6. http://”URLPROYECTO”.cl/disponibilidades/relaciones
    \begin{enumerate}
        \item a. Tipo de petición:
            \begin{itemize}
               \item i. “get”
            \end{itemize}
        \item b. Parámetros a ingresar, el tipo de datos y restricciones
            \begin{enumerate}
                \item i. Sin parámetros de entrada
            \end{enumerate}
    \end{enumerate}

    \item 7. http://”URLPROYECTO”.cl/disponibilidades/relaciones
    \begin{enumerate}
        \item a. Tipo de petición:
            \begin{itemize}
               \item i. “post”
            \end{itemize}
        \item b. Parámetros a ingresar, el tipo de datos y restricciones
            \begin{enumerate}
                \item i. id\_tipoe
                    \begin{enumerate}
                        \item 1. Number (entero)
                        \item 2. Requerido
                        \item 3. Debe existir en Tipo de Estado
                    \end{enumerate}
                \item ii. id\_dispo
                    \begin{enumerate}
                        \item 1. Number (entero)
                        \item 2. Requerido
                        \item 3. Debe existir en Disponibilidad
                        \item 4. La combinación id\_tipoe + id\_dispo no puede estar duplicada
                    \end{enumerate}
            \end{enumerate}
    \end{enumerate}

    \item 8. http://”URLPROYECTO”.cl/disponibilidades/relaciones/\{id\_tipoe\}/\{id\_dispo\}
    \begin{enumerate}
        \item a. Tipo de petición:
            \begin{itemize}
               \item i. “delete”
            \end{itemize}
        \item b. Parámetros a ingresar, el tipo de datos y restricciones
            \begin{enumerate}
                \item i. id\_tipoe (parámetro de ruta)
                    \begin{enumerate}
                        \item 1. Number (entero)
                        \item 2. Requerido
                    \end{enumerate}
                \item ii. id\_dispo (parámetro de ruta)
                    \begin{enumerate}
                        \item 1. Number (entero)
                        \item 2. Requerido
                    \end{enumerate}
            \end{enumerate}
        \item c. Nota: borrado físico
    \end{enumerate}
\end{enumerate}

\subsection{Marcas (/marcas)}

\begin{enumerate}
    \item 1. http://”URLPROYECTO”.cl/marcas/
    \begin{enumerate}
        \item a. Tipo de petición:
            \begin{itemize}
               \item i. “get”
            \end{itemize}
        \item b. Parámetros a ingresar, el tipo de datos y restricciones
            \begin{enumerate}
                \item i. Sin parámetros de entrada
            \end{enumerate}
    \end{enumerate}

    \item 2. http://”URLPROYECTO”.cl/marcas/\{id\_marca\}
    \begin{enumerate}
        \item a. Tipo de petición:
            \begin{itemize}
               \item i. “get”
            \end{itemize}
        \item b. Parámetros a ingresar, el tipo de datos y restricciones
            \begin{enumerate}
                \item i. id\_marca (parámetro de ruta)
                    \begin{enumerate}
                        \item 1. Number (entero)
                        \item 2. Requerido
                    \end{enumerate}
            \end{enumerate}
    \end{enumerate}

    \item 3. http://”URLPROYECTO”.cl/marcas/
    \begin{enumerate}
        \item a. Tipo de petición:
            \begin{itemize}
               \item i. “post”
            \end{itemize}
        \item b. Parámetros a ingresar, el tipo de datos y restricciones
            \begin{enumerate}
                \item i. nombre
                    \begin{enumerate}
                        \item 1. String
                        \item 2. Requerido
                        \item 3. Tamaño máximo de 100 caracteres
                    \end{enumerate}
            \end{enumerate}
    \end{enumerate}

    \item 4. http://”URLPROYECTO”.cl/marcas/\{id\_marca\}
    \begin{enumerate}
        \item a. Tipo de petición:
            \begin{itemize}
               \item i. “put”
            \end{itemize}
        \item b. Parámetros a ingresar, el tipo de datos y restricciones
            \begin{enumerate}
                \item i. id\_marca (parámetro de ruta)
                    \begin{enumerate}
                        \item 1. Number (entero)
                        \item 2. Requerido
                    \end{enumerate}
                \item ii. nombre
                    \begin{enumerate}
                        \item 1. String
                        \item 2. No requerido - Default null
                        \item 3. Tamaño máximo de 100 caracteres
                    \end{enumerate}
            \end{enumerate}
    \end{enumerate}

    \item 5. http://”URLPROYECTO”.cl/marcas/\{id\_marca\}
    \begin{enumerate}
        \item a. Tipo de petición:
            \begin{itemize}
               \item i. “delete”
            \end{itemize}
        \item b. Parámetros a ingresar, el tipo de datos y restricciones
            \begin{enumerate}
                \item i. id\_marca (parámetro de ruta)
                    \begin{enumerate}
                        \item 1. Number (entero)
                        \item 2. Requerido
                        \item 3. Bloqueado (400) si tiene Modelos asociados
                    \end{enumerate}
            \end{enumerate}
        \item c. Nota: borrado físico
    \end{enumerate}
\end{enumerate}

\subsection{Modelos (/modelos)}

\begin{enumerate}
    \item 1. http://”URLPROYECTO”.cl/modelos/
    \begin{enumerate}
        \item a. Tipo de petición:
            \begin{itemize}
               \item i. “get”
            \end{itemize}
        \item b. Parámetros a ingresar, el tipo de datos y restricciones
            \begin{enumerate}
                \item i. id\_marca (parámetro de query)
                    \begin{enumerate}
                        \item 1. Number (entero)
                        \item 2. No requerido - Default null
                    \end{enumerate}
            \end{enumerate}
    \end{enumerate}

    \item 2. http://”URLPROYECTO”.cl/modelos/\{id\_modelo\}
    \begin{enumerate}
        \item a. Tipo de petición:
            \begin{itemize}
               \item i. “get”
            \end{itemize}
        \item b. Parámetros a ingresar, el tipo de datos y restricciones
            \begin{enumerate}
                \item i. id\_modelo (parámetro de ruta)
                    \begin{enumerate}
                        \item 1. Number (entero)
                        \item 2. Requerido
                    \end{enumerate}
            \end{enumerate}
    \end{enumerate}

    \item 3. http://”URLPROYECTO”.cl/modelos/
    \begin{enumerate}
        \item a. Tipo de petición:
            \begin{itemize}
               \item i. “post”
            \end{itemize}
        \item b. Parámetros a ingresar, el tipo de datos y restricciones
            \begin{enumerate}
                \item i. nombre
                    \begin{enumerate}
                        \item 1. String
                        \item 2. Requerido
                        \item 3. Tamaño máximo de 100 caracteres
                    \end{enumerate}
                \item ii. id\_marca
                    \begin{enumerate}
                        \item 1. Number (entero)
                        \item 2. Requerido
                        \item 3. Debe existir en Marca
                    \end{enumerate}
            \end{enumerate}
    \end{enumerate}

    \item 4. http://”URLPROYECTO”.cl/modelos/\{id\_modelo\}
    \begin{enumerate}
        \item a. Tipo de petición:
            \begin{itemize}
               \item i. “put”
            \end{itemize}
        \item b. Parámetros a ingresar, el tipo de datos y restricciones
            \begin{enumerate}
                \item i. id\_modelo (parámetro de ruta)
                    \begin{enumerate}
                        \item 1. Number (entero)
                        \item 2. Requerido
                    \end{enumerate}
                \item ii. nombre
                    \begin{enumerate}
                        \item 1. String
                        \item 2. No requerido - Default null
                        \item 3. Tamaño máximo de 100 caracteres
                    \end{enumerate}
                \item iii. id\_marca
                    \begin{enumerate}
                        \item 1. Number (entero)
                        \item 2. No requerido - Default null
                        \item 3. Si se envía, se revalida su existencia
                    \end{enumerate}
            \end{enumerate}
    \end{enumerate}

    \item 5. http://”URLPROYECTO”.cl/modelos/\{id\_modelo\}
    \begin{enumerate}
        \item a. Tipo de petición:
            \begin{itemize}
               \item i. “delete”
            \end{itemize}
        \item b. Parámetros a ingresar, el tipo de datos y restricciones
            \begin{enumerate}
                \item i. id\_modelo (parámetro de ruta)
                    \begin{enumerate}
                        \item 1. Number (entero)
                        \item 2. Requerido
                    \end{enumerate}
            \end{enumerate}
        \item c. Nota: borrado físico, sin bloqueo por dependencias
    \end{enumerate}
\end{enumerate}

\subsection{Dashboard (/dashboard)}

\begin{enumerate}
    \item 1. http://”URLPROYECTO”.cl/dashboard/resumen
    \begin{enumerate}
        \item a. Tipo de petición:
            \begin{itemize}
               \item i. “get”
            \end{itemize}
        \item b. Parámetros a ingresar, el tipo de datos y restricciones
            \begin{enumerate}
                \item i. anio (parámetro de query)
                    \begin{enumerate}
                        \item 1. Number (entero)
                        \item 2. No requerido - Default: año calendario actual
                        \item 3. Rango válido: 2000 a 2100 (422 si está fuera de rango)
                        \item 4. Filtra todos los agregados de activos por \texttt{fecha\_recep} y todos los de licencia por \texttt{fecha\_asignacion}
                    \end{enumerate}
            \end{enumerate}
        \item c. Respuesta 200:
            \begin{itemize}
                \item i. anio: año efectivo de la consulta
                \item ii. activos: \{ periferico, equipo\_computo, adquiridos\_anio \} — todos filtrados por \texttt{fecha\_recep} del año
                \item iii. licencias\_totales: lista \{ proveedor, total \} de asignaciones del año (\texttt{fecha\_asignacion})
                \item iv. sistemas\_operativos: lista \{ familia, total \} de equipos con \texttt{fecha\_recep} del año
                \item v. activos\_por\_estado: lista \{ estado, total \} de activos con \texttt{fecha\_recep} del año
                \item vi. licencias\_por\_mes: \{ labels (Ene–Dic), series: [\{ proveedor, data[12] \}] \} del año
            \end{itemize}
        \item d. Nota: requiere token de autenticación
    \end{enumerate}
\end{enumerate}

\subsection{Solicitudes de Baja (/solicitudes-baja)}

\begin{enumerate}
    \item 1. http://”URLPROYECTO”.cl/solicitudes-baja/
    \begin{enumerate}
        \item a. Tipo de petición:
            \begin{itemize}
               \item i. “get”
            \end{itemize}
        \item b. Parámetros a ingresar, el tipo de datos y restricciones
            \begin{enumerate}
                \item i. estado\_sol (parámetro de query)
                    \begin{enumerate}
                        \item 1. String
                        \item 2. No requerido - Default null
                        \item 3. Valores esperados: "pendiente", "aprobada", "rechazada"
                    \end{enumerate}
                \item ii. id\_act (parámetro de query)
                    \begin{enumerate}
                        \item 1. Number (entero)
                        \item 2. No requerido - Default null
                    \end{enumerate}
            \end{enumerate}
    \end{enumerate}

    \item 2. http://”URLPROYECTO”.cl/solicitudes-baja/\{id\_solbaja\}
    \begin{enumerate}
        \item a. Tipo de petición:
            \begin{itemize}
               \item i. “get”
            \end{itemize}
        \item b. Parámetros a ingresar, el tipo de datos y restricciones
            \begin{enumerate}
                \item i. id\_solbaja (parámetro de ruta)
                    \begin{enumerate}
                        \item 1. Number (entero)
                        \item 2. Requerido
                    \end{enumerate}
            \end{enumerate}
    \end{enumerate}

    \item 3. http://”URLPROYECTO”.cl/solicitudes-baja/
    \begin{enumerate}
        \item a. Tipo de petición:
            \begin{itemize}
               \item i. “post”
            \end{itemize}
        \item b. Parámetros a ingresar, el tipo de datos y restricciones
            \begin{enumerate}
                \item i. motivo\_baja
                    \begin{enumerate}
                        \item 1. String
                        \item 2. Requerido
                        \item 3. Tamaño máximo de 200 caracteres
                    \end{enumerate}
                \item ii. id\_act
                    \begin{enumerate}
                        \item 1. Number (entero)
                        \item 2. Requerido
                        \item 3. Debe existir en Activo TI
                        \item 4. No puede tener ya una solicitud vigente distinta de "rechazada"
                    \end{enumerate}
                \item iii. id\_tipoe
                    \begin{enumerate}
                        \item 1. Number (entero)
                        \item 2. Requerido
                        \item 3. Corresponde a la nueva condición del activo: solo se permite "regular" o "malo" (se rechaza "bueno")
                    \end{enumerate}
            \end{enumerate}
        \item c. Nota: id\_emp\_sol se asigna automáticamente con el usuario autenticado. Solo requiere autenticación, no rol editor.
    \end{enumerate}

    \item 4. http://”URLPROYECTO”.cl/solicitudes-baja/\{id\_solbaja\}/resolver
    \begin{enumerate}
        \item a. Tipo de petición:
            \begin{itemize}
               \item i. “put”
            \end{itemize}
        \item b. Parámetros a ingresar, el tipo de datos y restricciones
            \begin{enumerate}
                \item i. id\_solbaja (parámetro de ruta)
                    \begin{enumerate}
                        \item 1. Number (entero)
                        \item 2. Requerido
                    \end{enumerate}
                \item ii. estado\_sol
                    \begin{enumerate}
                        \item 1. String
                        \item 2. Requerido
                        \item 3. Debe cumplir el patrón "^(aprobada|rechazada)\$"
                    \end{enumerate}
                \item iii. ref\_aprob
                    \begin{enumerate}
                        \item 1. String
                        \item 2. No requerido - Default null
                        \item 3. Tamaño máximo de 200 caracteres
                    \end{enumerate}
                \item iv. id\_emp\_aprob
                    \begin{enumerate}
                        \item 1. Number (entero)
                        \item 2. No requerido - Default null
                        \item 3. Debe existir en Empleado
                    \end{enumerate}
            \end{enumerate}
        \item c. Nota: solo aplicable si la solicitud está en estado "pendiente"; fecha\_resolucion se asigna automáticamente
    \end{enumerate}

    \item 5. http://”URLPROYECTO”.cl/solicitudes-baja/\{id\_solbaja\}/borrado
    \begin{enumerate}
        \item a. Tipo de petición:
            \begin{itemize}
               \item i. “put”
            \end{itemize}
        \item b. Parámetros a ingresar, el tipo de datos y restricciones
            \begin{enumerate}
                \item i. id\_solbaja (parámetro de ruta)
                    \begin{enumerate}
                        \item 1. Number (entero)
                        \item 2. Requerido
                    \end{enumerate}
                \item ii. ruta\_acta\_baja
                    \begin{enumerate}
                        \item 1. String
                        \item 2. No requerido - Default null
                        \item 3. Tamaño máximo de 255 caracteres
                    \end{enumerate}
            \end{enumerate}
        \item c. Nota: solo aplicable si la solicitud está "aprobada" y sin fecha\_borrado registrada; actualiza automáticamente id\_dispo del activo a "de\_baja", fecha\_borrado e id\_emp\_borrado
    \end{enumerate}
\end{enumerate}
